This chapter is the blueprint's appendix.  It records Lemma~A.1 of the paper
(\texttt{lem:semigroup-spectral-growth}, in its Appendix~A, \texttt{sec:semigroup-growth}) ---
Lemma~\ref{lem:semigroup-spectral-growth} here, the analytic input that turns the spectrum of
$L_{\eps,u}$ into a statement about the growth rate $\gamma(u,\eps)$ --- together with the
facts about $\gamma(u,\eps)$ that §1.1 of the paper states with it (the limit
\texttt{eq:semigroup-growth-rate} exists, and the spectral supremum is attained), and the Lean
chain that proves them.  The other nodes of this chapter are the formalisation's own: the paper
states none of them, and they are compared with the paper only through
Lemma~\ref{lem:semigroup-spectral-growth}, the growth-rate clause of Theorem~\ref{thm:main}, the
energy identity \texttt{eq:energy-competition} and the phrase ``its solution'' of
Theorem~\ref{thm:main}.  The particle flow, its entropy, the ideal induction group and the
vanishing-diffusivity limit are Section~4 and Remark~1.7 of the paper and have their own
chapter.

The paper proves its Lemma~A.1 on $\Lp^2_\sigma$ directly: $\eps\Delta$ with domain
$\Hs^2\cap\Lp^2_\sigma$ generates an analytic semigroup there, $C_uB=\curl(u\times B)$ is
infinitesimally bounded relative to it, so $L_{\eps,u}$ generates an analytic semigroup by
the perturbation theorem [Engel--Nagel, Theorem III.2.10]; the graph norm is equivalent to
the $\Hs^2$ norm, so the resolvent is compact; analyticity makes the semigroup immediately
compact, which gives the spectral mapping theorem and the equality of the growth bound with
the spectral bound; and $0\in\sigma(L_{\eps,u})$, together with the finiteness of the
spectrum in any right half-plane, makes the spectral supremum a maximum.  The formalisation
proves every one of these statements (Lemma~\ref{lem:semigroup-spectral-growth},
\texttt{SemigroupSpectralGrowth}), along a different route.  Mathlib has no
$\C_0$-semigroup theory, so the development builds what it needs from scratch.
\begin{itemize}
  \item \emph{Generation} (Sections~\ref{sec:dynamics:hille-yosida}
    to~\ref{sec:dynamics:solenoidal}).  A resolvent estimate $\norm{R(r)}\leq(r-M)^{-1}$ on a
    real ray (Definition~\ref{def:dynamics:hille-yosida-ray}) drives a Yosida approximation
    scheme whose limit is a strongly continuous semigroup
    (Theorem~\ref{thm:dynamics:hille-yosida}).  Feeding it the dissipativity of the reduced
    operator $A_{\eps,K}$ and of its formal adjoint
    (Lemma~\ref{lem:dynamics:reduced-adjoint},
    Proposition~\ref{prop:dynamics:reduced-half-plane-resolvent}) gives each axial sector a
    Hille--Yosida ray (Theorem~\ref{thm:dynamics:sector-semigroup}); summing the sector
    resolvents over the axial Fourier index (Lemmas~\ref{lem:dynamics:axial-hilbert-sum}
    and~\ref{lem:dynamics:uniform-sector-estimate},
    Proposition~\ref{prop:dynamics:product-resolvent}) gives a ray for, and hence generates
    the semigroup of, the all-frequency auxiliary operator $\mathcal A_\eps$
    (\texttt{productInductionOperator u $\eps$}, domain $\Hs^2(\T^3;\CC^3)$) on the whole of
    $\Lp^2(\T^3;\CC^3)$ (Theorem~\ref{thm:dynamics:product-semigroup}); restricting it to the
    closed solenoidal subspace (Lemma~\ref{lem:dynamics:solenoidal-subspace},
    Theorem~\ref{thm:dynamics:solenoidal-semigroup}) puts it on the paper's state space
    $\Lp^2_\sigma(\T^3)$.
  \item \emph{The generator} (Section~\ref{sec:dynamics:generator}).  The generator of a
    $\C_0$-semigroup and the spectrum of an unbounded operator are defined
    (Definition~\ref{def:dynamics:generator}); the generator of the Hille--Yosida semigroup is
    the operator it was built from (Lemma~\ref{lem:dynamics:hille-yosida-generator}); and the
    generator of the semigroup on $\Lp^2_\sigma$ is the paper's $L_{\eps,u}$ with domain
    $\Hs^2\cap\Lp^2_\sigma$, which determines the semigroup
    (Theorem~\ref{thm:dynamics:solenoidal-generator}).  So the semigroup whose growth bound is
    the Lean value of $\gamma(u,\eps)$ (Definition~\ref{def:main:growth-rate},
    \texttt{FastDynamo.gamma}) is the paper's $\e^{tL_{\eps,u}}$.
  \item \emph{Analyticity, compactness, spectral mapping} (Section~\ref{sec:dynamics:analytic}).
    A sector condition on the numerical range of $\mathcal A_\eps$
    (Lemma~\ref{lem:dynamics:sectorial}) passes to the Yosida approximants in complex time and
    makes the semigroup analytic of angle $\pi/4$
    (Theorem~\ref{thm:dynamics:analytic-generation}); the resolvent factors through the compact
    embedding $\Hs^2\hookrightarrow\Lp^2$ (Lemma~\ref{lem:dynamics:compact-embedding}); an
    analytic semigroup with compact resolvent is immediately compact
    (Lemma~\ref{lem:dynamics:immediate-compactness}); and for such semigroups a direct
    point-spectrum argument gives the spectral mapping theorem and the equality of the growth
    bound with the spectral bound, attained at an eigenvalue
    (Theorem~\ref{thm:dynamics:spectral-mapping}).  Mean conservation gives
    $0\in\sigma(L_{\eps,u})$ and $\norm{\e^{tL_{\eps,u}}}\geq1$
    (Lemma~\ref{lem:dynamics:mean-conservation}), and the growth bound of every
    $\C_0$-semigroup is a limit (Lemma~\ref{lem:dynamics:growth-bound-limit}).
    Theorem~\ref{thm:dynamics:induction-spectrum} collects the instances at the induction
    semigroup.
  \item \emph{Solutions} (Section~\ref{sec:dynamics:solutions}).  The energy identity
    \texttt{eq:energy-competition} (Lemma~\ref{lem:dynamics:energy-identity}), and uniqueness
    of strong solutions of the curl-form induction equation, which identifies every solution
    with divergence-free datum with the semigroup orbit
    (Lemma~\ref{lem:dynamics:physical-uniqueness}).
\end{itemize}

The growth-rate clause of Theorem~\ref{thm:main} needs only part of this: the paper bounds
$\norm{\e^{tL_{\eps,u}}}$ from below by $\norm{B^\eps(t)}_2/\norm{B_{\rm in}^\eps}_2$, and Lean
bounds the growth bound from below by the real part of an eigenvalue whose eigenfunction lies in
$\Lp^2_\sigma$ (Theorem~\ref{thm:main-growth-rate}); analyticity, compactness and spectral
mapping are not used there.  Both the paper and the Lean clause (\texttt{gamma\_ge\_gamma0})
state the growth-rate clause for $0<\eps\leq\eps_0$.

Besides the semigroup chain, the first section records the Lean-side nodes: the
$\C_0$-semigroup interface (Definition~\ref{def:dynamics:semigroup-interface}); the a priori
energy bound $\norm{B(t)}_2\leq\e^{M_ut}\norm{B(0)}_2$ for forward strong solutions of the
reduced operator and the uniqueness of those solutions
(Lemmas~\ref{lem:dynamics:energy-bound} and~\ref{lem:dynamics:strong-solution-uniqueness});
the exponential growth bound of an \emph{arbitrary} strongly continuous semigroup, its two
elementary norm bounds and its description as a limit
(Definition~\ref{def:dynamics:growth-bound},
Lemmas~\ref{lem:dynamics:growth-bound-norm-bounds} and~\ref{lem:dynamics:growth-bound-limit});
the exact exponential law of an explicit eigen-trajectory
(Lemma~\ref{lem:dynamics:eigen-trajectory-law}); and the growth floor of the initial-value
problem on the operator domain, sandwiched at the explicit witnesses by
$\gamma_0\leq\gamma_{\rm str}(A_{\eps,K_\eps})\leq M_u$
(Lemma~\ref{lem:dynamics:strong-solution-growth-floor}).  Throughout,
$M_u:=\norm{\nabla_\perp u}_\infty$ is the $\Lp^\infty$ norm of the stretching matrix field,
independent of $\eps$ and of the axial index.  The map \texttt{docs/MODULE\_MAP.md} contains
the same accounting on the Lean side.

\section{The growth rate and the analytic semigroup}

\begin{lemma}[Semigroup generation and spectral growth]
  \label{lem:semigroup-spectral-growth}
  \lean{FastDynamo.SemigroupSpectralGrowth, FastDynamo.semigroupSpectralGrowth_holds}
  \leanok
  \uses{def:main:induction-operator, def:main:growth-rate, def:dynamics:generator,
    def:dynamics:analytic-semigroup}
  Let $u\in\W^{1,\infty}(\T^3;\R^3) $ be divergence-free and let $\eps>0 $.  Then the
  following hold.
  \begin{enumerate}
    \item The operator $L_{\eps,u} $, with domain $\Hs^2\cap\Lp^2_\sigma $, generates an
      analytic $\C_0 $-semigroup on $\Lp^2_\sigma $.
    \item The resolvent of $L_{\eps,u} $ is compact, $\e^{tL_{\eps,u}} $ is compact for
      every $t>0 $, and
      \begin{equation}
        \sigma(\e^{tL_{\eps,u}})\setminus\{0\}
        =\left\{\e^{t\lambda}:\lambda\in\sigma(L_{\eps,u})\right\},
        \qquad t>0.
        \label{eq:dynamics:compact-semigroup-spectral-mapping}
      \end{equation}
    \item Furthermore,
      \begin{equation}
        \gamma(u,\eps)
        =\max\{\re\lambda:\lambda\in\sigma(L_{\eps,u})\}.
        \label{eq:dynamics:growth-rate-spectral-bound}
      \end{equation}
  \end{enumerate}
\end{lemma}

The paper states the lemma in two paragraphs, the second with the two displays; the
enumeration into items (1)--(3) is the blueprint's, for reference.  The Lean statement is the
structure \texttt{SemigroupSpectralGrowth u hu h$\eps$}, proved for every divergence-free
transverse velocity and every $\eps>0$ by \texttt{semigroupSpectralGrowth\_holds}
(\texttt{FastDynamo/Dynamics/Assembly/SemigroupSpectralGrowth.lean}).  Its fields, in order:
\texttt{generator\_eq}, \texttt{generator\_curlForm}, \texttt{generator\_unique} and
\texttt{analytic} (item~(1)); \texttt{resolventSet\_nonempty},
\texttt{isCompactOperator\_resolvent}, \texttt{isCompactOperator\_semigroup} and
\texttt{spectralMapping} (item~(2)); \texttt{gamma\_tendsto} (the limit
\texttt{eq:semigroup-growth-rate} exists and is a real number, as §1.1 of the paper states);
\texttt{gamma\_eq\_iSup} ($\gamma(u,\eps)$ is the supremum);
\texttt{zero\_mem\_spectrum}; and \texttt{exists\_eigenvalue\_re\_eq\_gamma} (the supremum is
attained, at an eigenvalue with an eigenvector, as §1.1 of the paper states);
\texttt{gamma\_eq\_iSup} and \texttt{exists\_eigenvalue\_re\_eq\_gamma} together are
\eqref{eq:dynamics:growth-rate-spectral-bound}.  It differs from the paper's
statement as follows.
\begin{itemize}
  \item \emph{Model.}  The torus is the product $\T^2\times\T$, the state space is the
    Fourier-defined $\Lp^2_\sigma$ of that model (\texttt{productSolenoidalSubspace}), and the
    velocity is axially independent (\texttt{TransverseRealW1Inf}), as everywhere in this
    blueprint.  The clauses are stated for the semigroup \texttt{solenoidalInductionSemigroup}
    on the product model; \texttt{gamma} is the growth bound of its isometric transport to the
    literal three-torus, which has the same operator norms and the same growth bound
    (\texttt{growthBound\_literalInductionSemigroup}).
  \item \emph{The operator.}  The generator is \texttt{solenoidalInductionOperator}: domain
    the divergence-free $\Hs^2$ fields, values the auxiliary action $\mathcal A_\eps B$ of
    Definition~\ref{def:main:induction-operator}, which lies in $\Lp^2_\sigma$.  On that
    domain the action is the curl form: its image in the Fourier model of $\Hs^{-2}$ is
    $\curl(u\times B)+\eps\Delta B$ (\texttt{generator\_curlForm}), and the embedding of
    $\Lp^2$ into that model is injective.  This is an equivalent formulation of the paper's
    operator.
  \item \emph{Analyticity} is \texttt{IsAnalyticSemigroup} with the explicit angle $\pi/4$
    (Definition~\ref{def:dynamics:analytic-semigroup}); the paper names no angle.  The
    spectrum is \texttt{LinearPMap.spectrum}, the complement of the set of $z$ at which
    $z-L$ has a bounded two-sided inverse (Definition~\ref{def:dynamics:generator}).
  \item \emph{Additional clauses.}  \texttt{generator\_unique} (every $\C_0$-semigroup on
    $\Lp^2_\sigma$ with this generator is this one), \texttt{resolventSet\_nonempty} and
    \texttt{zero\_mem\_spectrum} are not in the paper's statement; the last is a step of the
    paper's proof.  Attainment is stated with an eigenvector.
\end{itemize}
The route differs from the paper's at every step; the proof below gives both.

\begin{proof}
  \leanok
  \uses{def:main:induction-operator, def:main:growth-rate,
    thm:dynamics:product-semigroup, thm:dynamics:solenoidal-semigroup,
    thm:dynamics:solenoidal-generator, thm:dynamics:induction-spectrum,
    lem:dynamics:growth-bound-limit, lem:dynamics:mean-conservation}
  \emph{The paper's proof.}  Work on $\Lp^2_\sigma $ throughout.  The operator $\eps\Delta $
  with domain $\Hs^2\cap\Lp^2_\sigma $ generates an analytic semigroup there by its Fourier
  representation, and $C_uB=\curl(u\times B) $ belongs to $\Lp^2_\sigma $ for every
  $B\in\Hs^2\cap\Lp^2_\sigma $.  For every $\eta>0 $, interpolation gives
  $\norm{C_uB}_2\leq C\norm u_{\W^{1,\infty}}\norm B_{\Hs^1}\leq\eta\norm{\Delta B}_2+C_\eta\norm B_2 $,
  so $C_u $ is infinitesimally bounded relative to $\eps\Delta $ on $\Lp^2_\sigma $, and the
  perturbation theorem for analytic semigroups [Engel--Nagel, Theorem III.2.10] gives item
  (1) with the stated domain.  By the same relative bound and the elliptic estimate
  $\norm B_{\Hs^2}\leq C(\norm{\Delta B}_2+\norm B_2) $ the graph norm of $L_{\eps,u} $ is
  equivalent to the $\Hs^2 $-norm, so the compact embedding
  $\Hs^2(\T^3)\hookrightarrow\Lp^2(\T^3) $ gives a compact resolvent; analyticity plus
  [Engel--Nagel, Theorem II.4.29] makes the semigroup immediately compact, and
  [Engel--Nagel, Corollaries IV.3.11 and IV.3.12] give
  \eqref{eq:dynamics:compact-semigroup-spectral-mapping} and the equality of the growth
  bound with the spectral bound, which is
  \eqref{eq:dynamics:growth-rate-spectral-bound}.  For attainment: every
  $B\in D(L_{\eps,u}) $ has $\int_{\T^3}L_{\eps,u}B\dd x=0 $, and $\Lp^2_\sigma $ contains
  non-zero constant fields, so $L_{\eps,u} $ is not surjective and $0\in\sigma(L_{\eps,u}) $.
  By [Engel--Nagel, Corollary V.3.2] only finitely many spectral points lie in any fixed
  right half-plane; applied to $\{\re z\geq0\} $, this makes the spectral supremum a maximum.
  The paper's proof notes in addition that the limit in \texttt{eq:semigroup-growth-rate} is
  this growth bound.

  \emph{The Lean proof.}  Item~(1): the semigroup is generated on all of
  $\Lp^2(\T^3;\CC^3)$ by the auxiliary operator through the Hille--Yosida theorem
  (Theorem~\ref{thm:dynamics:product-semigroup}) and restricted to the invariant closed
  subspace $\Lp^2_\sigma$ (Theorem~\ref{thm:dynamics:solenoidal-semigroup}); its generator is
  identified with $L_{\eps,u}$ on $\Hs^2\cap\Lp^2_\sigma$ and determines it
  (Theorem~\ref{thm:dynamics:solenoidal-generator}); analyticity comes from a sector condition
  on the numerical range through the Yosida approximants in complex time, not from
  [Engel--Nagel, Theorem III.2.10] (Theorem~\ref{thm:dynamics:analytic-generation}).
  Item~(2): the resolvent factors through the compact embedding, as in the paper;
  compactness of $\e^{tL_{\eps,u}}$ for $t>0$ is proved from analyticity directly
  (Lemma~\ref{lem:dynamics:immediate-compactness}), where the paper cites
  [Engel--Nagel, Theorem II.4.29]; and the spectral mapping theorem is proved by a direct
  point-spectrum argument (Theorem~\ref{thm:dynamics:spectral-mapping}), where the paper cites
  [Engel--Nagel, Corollaries IV.3.11 and IV.3.12].  Item~(3): the same point-spectrum argument,
  with Gelfand's formula for $\e^{L_{\eps,u}}$ in place of [Engel--Nagel, Corollary V.3.2],
  gives the growth bound as the spectral bound, attained at an eigenvalue; the spectrum is
  non-empty because $0\in\sigma(L_{\eps,u})$, which is the paper's mean argument
  (Lemma~\ref{lem:dynamics:mean-conservation}).  The limit is
  Lemma~\ref{lem:dynamics:growth-bound-limit}, and it is a real number because
  $\norm{\e^{tL_{\eps,u}}}\geq1$ (Lemma~\ref{lem:dynamics:mean-conservation}).  The instances
  at the induction semigroup are Theorem~\ref{thm:dynamics:induction-spectrum}.
\end{proof}

\begin{definition}[The Lean $\C_0$-semigroup interface]
  \label{def:dynamics:semigroup-interface}
  \lean{FastDynamo.StronglyContinuousSemigroup, FastDynamo.DomainOperator.IsForwardStrongSolution, FastDynamo.DomainOperator.ForwardStrongSolutionsUnique, FastDynamo.DomainOperatorSemigroupRealization}
  \leanok
  \uses{def:main:induction-operator}
  For a complex normed space $X $, a \emph{strongly continuous semigroup} is a family
  $S(t)\in\mathcal L(X) $, $t\geq0 $, with $S(0)=\id $,
  $S(s+t)=S(s)S(t) $, and $t\mapsto S(t)x $ continuous for every $x\in X $.  For an
  operator $A $ presented as a bounded map out of a Banach domain $D $ with continuous
  injection $\iota:D\to X $, a curve $B:\R\to D $ is a \emph{forward strong solution}
  when $t\mapsto\iota(B(t)) $ has one-sided derivative $A(B(t)) $ at every $t\geq0 $, and
  $A $ has \emph{unique forward strong solutions} when two forward strong solutions with
  the same initial ambient state agree for all $t\geq0 $.  A \emph{realization} of $S $ by
  $A $ is the data of a domain-valued orbit $B\mapsto\Theta(B,\cdot) $ with
  $\Theta(B,0)=B $, each $\Theta(B,\cdot) $ a forward strong solution, and
  $S(t)\iota(B)=\iota(\Theta(B,t)) $ for every $t\geq0 $.
\end{definition}

In Lean the four items are \texttt{StronglyContinuousSemigroup} (a structure bundling the
three semigroup axioms, with time in \texttt{NNReal}),
\texttt{DomainOperator.IsForwardStrongSolution},
\texttt{DomainOperator.ForwardStrongSolutionsUnique} and
\texttt{DomainOperatorSemigroupRealization}, all in
\texttt{FastDynamo/Dynamics/StronglyContinuousSemigroup.lean}.  The realization structure
packages the \emph{conclusion} of item~(1) of Lemma~\ref{lem:semigroup-spectral-growth} as a
hypothesis, and \emph{both} of the hypotheses this interface isolates are discharged for
the operators of this development, so the theorems quantified over a realization are not
vacuous:
Lemma~\ref{lem:dynamics:strong-solution-uniqueness} and
Lemma~\ref{lem:dynamics:product-dissipative} inhabit
\texttt{ForwardStrongSolutionsUnique} for $A_{\eps,K} $ and for
$\mathcal A_\eps $, at every divergence-free $u $ and $\eps\geq0 $, and \texttt{HilleYosidaRay.realization}
(Theorem~\ref{thm:dynamics:hille-yosida}) and
\texttt{productInductionRealization} (Theorem~\ref{thm:dynamics:product-semigroup})
inhabit \texttt{DomainOperatorSemigroupRealization}.  One reading of the interface should
be resisted: a realization records that the semigroup \emph{solves} the Cauchy
problem, not that it is the only family doing so, and it is forward uniqueness, not the
structure, that supplies the definite article
(Lemma~\ref{lem:dynamics:semigroup-unique}).

\begin{lemma}[A priori energy bound for forward strong solutions]
  \label{lem:dynamics:energy-bound}
  \lean{FastDynamo.DomainOperator.norm_embedding_le_exp_of_isForwardStrongSolution, FastDynamo.reducedInductionOperator_norm_embedding_le_exp, FastDynamo.re_inner_reducedInductionAction_le}
  \leanok
  \uses{def:main:induction-operator}
  Let $A $ be an operator with bundled domain $D $ and continuous injection $\iota:D\to X $
  into a complex inner-product space, and suppose that the numerical range of $A $ lies in
  a half-plane,
  \[
    \re\inner{\iota(d)}{A d}\leq M\norm{\iota(d)}^2
    \qquad(d\in D).
  \]
  Then every forward strong solution $B $ of $A $ obeys
  \[
    \norm{\iota(B(t))}\leq\e^{Mt}\norm{\iota(B(0))}
    \qquad(t\geq0).
  \]
  The hypothesis holds for the reduced fixed-sector operator $A_{\eps,K} $ of
  \eqref{eq:main:torus-sector} with $M=M_u:=\norm{\nabla_\perp u}_\infty $, for every
  divergence-free transverse $u\in\W^{1,\infty} $, every $\eps\geq0 $ and every $K\in\Z $;
  hence every forward strong solution of $A_{\eps,K} $ satisfies
  $\norm{B(t)}_2\leq\e^{M_ut}\norm{B(0)}_2 $ for $t\geq0 $, with a ceiling independent of
  the diffusivity and of the axial index.
\end{lemma}

This is the operator-level form of the paper's energy competition
(\texttt{eq:energy-competition}): stretching contributes at most
$M_u\norm B_2^2 $,
diffusion contributes a nonpositive term and transport contributes nothing.  In Lean the
abstract bound is
\texttt{DomainOperator.norm\_embedding\_le\_exp\_of\_isForwardStrongSolution} and its
instance for $A_{\eps,K} $ is
\texttt{reducedInductionOperator\_norm\_embedding\_le\_exp}, both in
\texttt{FastDynamo/Dynamics/StrongSolutionUniqueness.lean}; the half-plane hypothesis is
discharged there by \texttt{re\_inner\_reducedInductionAction\_le}, which combines the
exact energy identity and the stretching estimate of
\texttt{FastDynamo/Global/ReducedInductionEnergyIdentity.lean}.  The constant $M_u $ is the
$\Lp^\infty $ operator norm of the stretching matrix field
$b\mapsto(b_\perp\cdot\nabla_\perp)u $, and it does not depend on $\eps $ or $K $.

\begin{proof}
  \leanok
  \uses{def:main:induction-operator}
  Along a forward strong solution the squared ambient energy
  $\varphi(t)=\norm{\iota(B(t))}^2 $ is differentiable at every $t>0 $ with
  $\varphi'(t)=2\re\inner{\iota(B(t))}{A(B(t))}\leq2M\varphi(t) $, so
  $t\mapsto\e^{-2Mt}\varphi(t) $ has nonpositive derivative on $(0,T) $ and is antitone on
  $[0,T] $; evaluating at the endpoints and taking square roots gives the bound.  (This is
  Gronwall's lemma in exponential-multiplier form; Lean carries it out by hand, with
  mean-value theory only, rather than through \texttt{Mathlib.Analysis.ODE.Gronwall}.)
  For $A_{\eps,K} $ the
  half-plane hypothesis is the exact identity
  \[
    \re\inner{b}{A_{\eps,K}b}
    =-\eps\bigl(\norm{\nabla_\perp b}_2^2+(2\pi K)^2\norm b_2^2\bigr)
    +\re\inner{b}{(b_\perp\cdot\nabla_\perp)u},
  \]
  in which the transport terms are skew and drop out, together with the stretching bound
  $\re\inner{b}{(b_\perp\cdot\nabla_\perp)u}\leq M_u\norm b_2^2 $ and $\eps\geq0 $.
\end{proof}

\begin{lemma}[Uniqueness of forward strong solutions]
  \label{lem:dynamics:strong-solution-uniqueness}
  \lean{FastDynamo.DomainOperator.forwardStrongSolutionsUnique_of_re_inner_le, FastDynamo.reducedInductionOperator_forwardStrongSolutionsUnique}
  \leanok
  \uses{lem:dynamics:energy-bound}
  With $A $, $\iota $ and $M $ as in Lemma~\ref{lem:dynamics:energy-bound}, two forward
  strong solutions of $A $ with the same initial ambient state agree in $X $ at every
  $t\geq0 $.  In particular, for every divergence-free transverse $u\in\W^{1,\infty} $,
  every $\eps\geq0 $ and every $K\in\Z $, the reduced operator $A_{\eps,K} $ has unique
  forward strong solutions, with no further hypothesis.
\end{lemma}

This is the uniqueness half of an energy well-posedness statement, for the
fixed-sector reduced operator; existence is not formalised.  The paper states no such
proposition, and Theorem~\ref{thm:main} does not need one.  The Lean statement is weaker
than a full energy well-posedness proposition in three ways -- the velocity is
autonomous, the data lie in the bundled $\Hs^2 $ domain rather than in $\Lp^2_\sigma $, and
the solution concept is the strong one of Definition~\ref{def:dynamics:semigroup-interface}
rather than the Lions--Magenes energy one -- and it holds for every divergence-free $u $, every
$\eps\geq0 $ and every $K $.  In Lean:
\texttt{DomainOperator.forwardStrongSolutionsUnique\_of\_re\_inner\_le} and
\texttt{reducedInductionOperator\_forwardStrongSolutionsUnique}
(\texttt{FastDynamo/Dynamics/StrongSolutionUniqueness.lean}).  It inhabits one of the two hypotheses isolated by
Definition~\ref{def:dynamics:semigroup-interface}.

\begin{proof}
  \leanok
  \uses{lem:dynamics:energy-bound}
  The difference of two forward strong solutions is again a forward strong solution, by
  linearity of $\iota $ and of $A $, and its initial ambient state vanishes.  The bound of
  Lemma~\ref{lem:dynamics:energy-bound} applied to the difference gives
  $\norm{\iota(B_1(t)-B_2(t))}\leq\e^{Mt}\cdot0=0 $ for every $t\geq0 $.
\end{proof}

\begin{lemma}[Exact exponential law of an eigen-trajectory]
  \label{lem:dynamics:eigen-trajectory-law}
  \lean{FastDynamo.DomainOperator.isStrongSolution_eigenTrajectory, FastDynamo.DomainOperator.norm_embedding_eigenTrajectory}
  \leanok
  \uses{def:main:induction-operator}
  Let $A $ be an operator with bundled Banach domain $D $ and continuous injection
  $\iota:D\to X $, and let $\lambda\in\CC $, $0\neq B\in D $ satisfy $AB=\lambda\iota(B) $.
  Then the curve $t\mapsto\e^{t\lambda}B $ is a strong solution of
  $\frac{\dd}{\dd t}\iota(B(t))=A(B(t)) $ on all of $\R $, and
  \[
    \norm{\iota(\e^{t\lambda}B)}
    =\e^{t\re\lambda}\norm{\iota(B)}
    \qquad(t\in\R).
  \]
\end{lemma}

In Lean the curve is \texttt{FastDynamo.DomainOperator.eigenTrajectory}, and the two
displayed facts are \texttt{isStrongSolution\_eigenTrajectory}
(\texttt{FastDynamo/Operator/Evolution.lean}) and
\texttt{norm\_embedding\_eigenTrajectory}
(\texttt{FastDynamo/Finish/GrowingModeEvolution.lean}); the same module derives that the
norm tends to infinity when $\re\lambda>0 $ and crosses every threshold at a nonnegative
time.

This is \emph{not} the growth bound of Lemma~\ref{lem:semigroup-spectral-growth}, and the
distinction is the whole content of this chapter.  The lemma describes a single
trajectory, written down by hand by scalar exponentiation; it asserts nothing about the
existence of an evolution operator, and the quantity it controls is
$\norm{\iota(\e^{t\lambda}B)} $ for one datum, not
$\norm{\e^{tL_{\eps,u}}}_{\Lp^2_\sigma\to\Lp^2_\sigma} $.  To convert it into a lower
bound for $\gamma(u,\eps) $ one needs a semigroup to exist and this curve to be its orbit,
which is exactly the realization hypothesis of
Definition~\ref{def:dynamics:semigroup-interface}; that hypothesis is a theorem
(Theorem~\ref{thm:dynamics:product-semigroup}), and the conversion is carried out in
Theorem~\ref{thm:main-growth-rate}.  It is this lemma, at the level of the physical
real-valued field, that supplies the exact norm law in Theorem~\ref{thm:main}; the
growth-rate clause is a separate statement, proved separately and by a different route.

\begin{proof}
  \leanok
  \uses{def:main:induction-operator}
  Differentiating $s\mapsto\e^{s\lambda} $ in the real variable $s $ and applying the
  bounded map $\iota $ gives
  $\frac{\dd}{\dd t}\iota(\e^{t\lambda}B)=\lambda\,\iota(\e^{t\lambda}B)$, which equals
  $A(\e^{t\lambda}B) $ by the eigenvalue equation and linearity.  For the norm,
  $\norm{\iota(\e^{t\lambda}B)}=\abs{\e^{t\lambda}}\norm{\iota(B)}=\e^{t\re\lambda}\norm{\iota(B)} $.
  Neither step needs a hypothesis or an analytic input.
\end{proof}

\begin{definition}[The growth bound of a strongly continuous semigroup]
  \label{def:dynamics:growth-bound}
  \lean{FastDynamo.StronglyContinuousSemigroup.growthBound}
  \leanok
  \uses{def:dynamics:semigroup-interface}
  For a strongly continuous semigroup $S $ on a complex normed space $X $, the
  \emph{exponential growth bound} is
  \[
    \omega_0(S)
    :=\limsup_{t\to\infty}\frac1t\log\norm{S(t)}
    \in[-\infty,+\infty],
  \]
  with the convention that a time at which $S(t)=0 $ contributes $-\infty $.
\end{definition}

In Lean this is \texttt{StronglyContinuousSemigroup.growthBound}
(\texttt{FastDynamo/Dynamics/GrowthBound.lean}), valued in \texttt{EReal}.  Two points of
the encoding matter.  First, it is defined as a $\limsup $, so that it is defined for every
semigroup with no preliminary work; Lemma~\ref{lem:dynamics:growth-bound-limit} proves that
for every $\C_0$-semigroup on a Banach space it is a limit, so it is the limit of
Definition~\ref{def:main:growth-rate}.  Second, the integrand is $-\infty $ rather than
Mathlib's junk value $\log0=0 $ at the times where $\norm{S(t)}=0 $, so the definition is
faithful on the nilpotent semigroups the interface permits.  This is the growth bound of an
arbitrary \emph{given} semigroup; applied to the semigroup of
Theorem~\ref{thm:dynamics:solenoidal-semigroup} it is the Lean value of $\gamma(u,\eps) $
(Definition~\ref{def:main:growth-rate}).

\begin{lemma}[Exponential majorants and minorants control the growth bound]
  \label{lem:dynamics:growth-bound-norm-bounds}
  \lean{FastDynamo.StronglyContinuousSemigroup.growthBound_le_of_norm_le, FastDynamo.StronglyContinuousSemigroup.le_growthBound_of_norm_ge}
  \leanok
  \uses{def:dynamics:growth-bound}
  Let $S $ be a strongly continuous semigroup on $X $ and let $\omega\in\R $.  If
  $\norm{S(t)}\leq C\e^{\omega t} $ for every $t\geq0 $ and some $C>0 $, then
  $\omega_0(S)\leq\omega $.  If $\e^{\omega t}\leq\norm{S(t)} $ for all large $t $, then
  $\omega\leq\omega_0(S) $.
\end{lemma}

\begin{proof}
  \leanok
  \uses{def:dynamics:growth-bound}
  For the first bound, at the times where $\norm{S(t)}=0 $ the integrand is $-\infty $ and
  there is nothing to prove; elsewhere take logarithms, divide by $t>0 $ and use
  monotonicity of the $\limsup $ against the majorant $t^{-1}\log C+\omega $, which tends
  to $\omega $.  For the second bound the hypothesis makes $\norm{S(t)} $ positive along
  $t\to\infty $, so the same computation shows that the integrand is eventually at least
  $\omega $, which is below the $\limsup $.
\end{proof}

\begin{lemma}[The growth bound is a limit]
  \label{lem:dynamics:growth-bound-limit}
  \lean{FastDynamo.StronglyContinuousSemigroup.growthExponent,
    FastDynamo.StronglyContinuousSemigroup.exists_norm_operator_le,
    FastDynamo.StronglyContinuousSemigroup.growthBound_eq_iInf,
    FastDynamo.StronglyContinuousSemigroup.tendsto_growthExponent,
    FastDynamo.StronglyContinuousSemigroup.growthBound_ne_top,
    FastDynamo.StronglyContinuousSemigroup.growthBound_ne_bot_iff,
    FastDynamo.StronglyContinuousSemigroup.exists_growthBound_eq_and_tendsto}
  \leanok
  \uses{def:dynamics:growth-bound}
  Let $S $ be a strongly continuous semigroup on a complex Banach space $X $.  Then
  \[
    \frac1t\log\norm{S(t)}\xrightarrow[t\to\infty]{}\omega_0(S)
    =\inf_{t>0}\frac1t\log\norm{S(t)}\in[-\infty,\infty),
  \]
  with $\log0=-\infty$.  The limit is finite exactly when there is $\omega\in\R$ with
  $\e^{\omega t}\leq\norm{S(t)}$ for every $t>0$, and then it is a real number $L$ with
  $t^{-1}\log\norm{S(t)}\to L$ in $\R$.  In particular the $\limsup$ of
  Definition~\ref{def:dynamics:growth-bound} is the limit of
  \eqref{eq:main:semigroup-growth-rate}.
\end{lemma}

This is the paper's justification of \texttt{eq:semigroup-growth-rate} (the sentence after that
display: subadditivity and local boundedness, [Engel--Nagel, Proposition IV.2.2]), for every
$\C_0$-semigroup, with no hypothesis beyond completeness of $X$.  In Lean:
\texttt{FastDynamo/Dynamics/GrowthBoundLimit.lean}.  \texttt{growthExponent} is the integrand of
\texttt{growthBound}, valued in \texttt{EReal} and equal to $\bot=-\infty$ where $S(t)=0$;
\texttt{growthBound\_eq\_iInf} and \texttt{tendsto\_growthExponent} are the display;
\texttt{growthBound\_ne\_bot\_iff} is the criterion for finiteness, and
\texttt{exists\_growthBound\_eq\_and\_tendsto} the real-valued form.  The limit lies in
$[-\infty,\infty)$, as in the paper; it is $-\infty$ for nilpotent semigroups and also for
semigroups whose norm decays faster than every exponential without vanishing.  Completeness of
$X$ enters through the Banach--Steinhaus theorem (\texttt{exists\_norm\_operator\_le}).

\begin{proof}
  \leanok
  \uses{def:dynamics:growth-bound}
  Fix $s>0$ and put $n=\lfloor t/s\rfloor$.  The semigroup law gives
  $\norm{S(t)}\leq\norm{S(t-ns)}\norm{S(s)}^n\leq K\norm{S(s)}^n$, where $K$ bounds
  $\norm{S(r)}$ for $0\leq r\leq s$ (strong continuity on a compact interval and the
  Banach--Steinhaus theorem).  Taking logarithms,
  $t^{-1}\log\norm{S(t)}\leq t^{-1}\log K+(n/t)\log\norm{S(s)}$, whose right side tends to
  $s^{-1}\log\norm{S(s)}$; if $S(s)=0$ then $S(t)=0$ for $t\geq s$.  Hence the $\limsup$ is
  at most every value $s^{-1}\log\norm{S(s)}$, so at most their infimum, which is at most the
  $\liminf$.  So $\limsup=\inf=\liminf$, and the limit exists.  It is finite exactly when the
  infimum is, which is the exponential minorant.
\end{proof}

\begin{lemma}[Spectral lower bound for the growth bound of a realizing semigroup]
  \label{lem:dynamics:conditional-growth-bound}
  \lean{FastDynamo.DomainOperatorSemigroupRealization.exp_re_le_operatorNorm, FastDynamo.DomainOperatorSemigroupRealization.eigenpair_orbit_and_growth, FastDynamo.DomainOperatorSemigroupRealization.re_le_growthBound_of_eigenpair}
  \leanok
  \uses{def:dynamics:growth-bound, def:dynamics:semigroup-interface, lem:dynamics:strong-solution-uniqueness, lem:dynamics:eigen-trajectory-law}
  Let $S $ be a strongly continuous semigroup on $X $, let $A $ be an operator with
  bundled domain $D $, assume a realization of $S $ by $A $ in the sense of
  Definition~\ref{def:dynamics:semigroup-interface}, and assume that $A $ has unique
  forward strong solutions.  Then for every eigenpair $AB=\lambda\iota(B) $ with
  $\iota(B)\neq0 $ and every $t\geq0 $,
  \[
    S(t)\iota(B)=\iota(\e^{t\lambda}B),
    \qquad
    \norm{S(t)\iota(B)}=\e^{t\re\lambda}\norm{\iota(B)},
    \qquad
    \e^{t\re\lambda}\leq\norm{S(t)},
  \]
  and consequently $\re\lambda\leq\omega_0(S) $.
\end{lemma}

The statement carries its own hypotheses, and in that form it is proved in Lean
(\texttt{FastDynamo/Dynamics/StronglyContinuousSemigroup.lean} for the three displayed
identities, \texttt{FastDynamo/Dynamics/GrowthBound.lean} for
\texttt{re\_le\_growthBound\_of\_eigenpair}).  It isolates the easy half of the growth
clause: the implication ``an eigenvalue with $\re\lambda\geq\gamma_0 $
$\Rightarrow$ $\norm{S(t)}\geq\e^{t\gamma_0} $, hence $\gamma(u,\eps)\geq\gamma_0 $''.  Both
of the hypotheses the lemma carries are theorems: forward
uniqueness for the reduced operator, for the all-frequency operator and for the curl form
(Lemmas~\ref{lem:dynamics:strong-solution-uniqueness},
\ref{lem:dynamics:product-dissipative} and~\ref{lem:dynamics:physical-uniqueness}), and the realization by the generation theorem
(Theorem~\ref{thm:dynamics:product-semigroup}).
The lemma is therefore not conditional in substance: it is the engine of the lower
bound of Lemma~\ref{lem:dynamics:hille-yosida-growth-sandwich} and, through it, of
Theorem~\ref{thm:main-growth-rate}.

\begin{proof}
  \leanok
  \uses{def:dynamics:growth-bound, lem:dynamics:growth-bound-norm-bounds, lem:dynamics:eigen-trajectory-law}
  The realization supplies a forward strong solution starting at $B $ whose embedding is
  the semigroup orbit; the eigen-trajectory of
  Lemma~\ref{lem:dynamics:eigen-trajectory-law} is a second forward strong solution with
  the same initial state.  Forward uniqueness identifies the two, which gives the first
  identity; the norm law of Lemma~\ref{lem:dynamics:eigen-trajectory-law} gives the
  second; and dividing $\norm{S(t)\iota(B)}\leq\norm{S(t)}\norm{\iota(B)} $ by
  $\norm{\iota(B)}>0 $ gives the third.  The third inequality holds for every $t\geq0 $, so
  the minorant half of Lemma~\ref{lem:dynamics:growth-bound-norm-bounds} with
  $\omega=\re\lambda $ gives $\re\lambda\leq\omega_0(S) $.  Forward uniqueness enters as a
  \emph{hypothesis} of this lemma, not as a lemma its proof invokes; it is discharged
  downstream at Lemma~\ref{lem:dynamics:strong-solution-uniqueness}, whose Lean module
  \texttt{Dynamics/StrongSolutionUniqueness.lean} imports
  \texttt{Dynamics/StronglyContinuousSemigroup.lean} and not the other way round, so this
  proof does not \texttt{\textbackslash uses} it.
\end{proof}

\begin{lemma}[The growth floor of the initial-value problem]
  \label{lem:dynamics:strong-solution-growth-floor}
  \lean{FastDynamo.DomainOperator.strongSolutionGrowthFloor, FastDynamo.DomainOperator.re_le_strongSolutionGrowthFloor_of_eigenpair, FastDynamo.DomainOperator.strongSolutionGrowthFloor_le_of_norm_le, FastDynamo.reducedInductionOperator_strongSolutionGrowthFloor_le, FastDynamo.mainGrowthFloor_le_strongSolutionGrowthFloor, FastDynamo.mainGrowthFloor_le_strongSolutionGrowthFloor_le}
  \leanok
  \uses{lem:dynamics:eigen-trajectory-law, lem:dynamics:energy-bound, thm:main-spectral, def:main:witnesses}
  For an operator $A $ with bundled domain $D $ and continuous injection $\iota:D\to X $
  put
  \[
    \gamma_{\rm str}(A)
    :=\sup_B\ \limsup_{t\to\infty}\frac1t
    \log\frac{\norm{\iota(B(t))}}{\norm{\iota(B(0))}}
    \in[-\infty,+\infty],
  \]
  the supremum being over all forward strong solutions $B $ of $A $ with
  $\iota(B(0))\neq0 $ (and $-\infty $ if there are none; a time at which $\iota(B(t)) $
  vanishes contributes $-\infty $ to the inner $\limsup $).  Then
  \[
    \re\lambda\leq\gamma_{\rm str}(A)
  \]
  for every eigenpair $AB=\lambda\iota(B) $ with $\iota(B)\neq0 $, and
  \[
    \gamma_{\rm str}(A)\leq M
  \]
  whenever every forward strong solution of $A $ obeys
  $\norm{\iota(B(t))}\leq\e^{Mt}\norm{\iota(B(0))} $ for $t\geq0 $.  Both halves apply to
  the reduced fixed-sector operator: with $M_u=\norm{\nabla_\perp u}_\infty $ as in
  Lemma~\ref{lem:dynamics:energy-bound} one has $\gamma_{\rm str}(A_{\eps,K})\leq M_u $
  for every divergence-free transverse $u\in\W^{1,\infty} $, every $\eps\geq0 $ and every
  $K\in\Z $; and at the explicit witnesses $u $, $\eps_0 $, $\gamma_0 $ of
  Definition~\ref{def:main:witnesses} the eigenpairs of Theorem~\ref{thm:main-spectral}
  give, for every $0<\eps\leq\eps_0 $, an axial index $K_\eps\neq0 $ with
  \[
    \gamma_0\leq\gamma_{\rm str}(A_{\eps,K_\eps})\leq M_u.
  \]
\end{lemma}

This has no hypothesis: no semigroup, no realization and no uniqueness hypothesis.  It is
\emph{not} the paper's $\gamma(u,\eps) $, and the Lean name
\texttt{strongSolutionGrowthFloor} is chosen so that it cannot be mistaken for it.  The
difference is the one this whole chapter is about: $\gamma_{\rm str} $ is the growth rate
of the initial-value problem on the operator's domain, a supremum over curves that happen
to solve the equation, whereas $\gamma(u,\eps) $ of Definition~\ref{def:main:growth-rate}
is a supremum over \emph{all} $\Lp^2_\sigma $ data taken inside an operator norm.  Given a
strongly continuous semigroup $S $ realizing $A $, the inequality
$\gamma_{\rm str}(A)\leq\omega_0(S) $ follows from
$\norm{S(t)x}\leq\norm{S(t)}\norm x $, and the converse from density of the domain
together with the locally uniform bound of Lemma~\ref{lem:dynamics:energy-bound}; neither
direction is proved.  In Lean:
\texttt{DomainOperator.strongSolutionGrowthFloor} and
\texttt{DomainOperator.re\_le\_strongSolutionGrowthFloor\_of\_eigenpair}
and the ceiling \texttt{DomainOperator.strongSolutionGrowthFloor\_le\_of\_norm\_le}
(\texttt{FastDynamo/Dynamics/GrowthBound.lean}).  The sandwich at the explicit witnesses is
\texttt{mainGrowthFloor\_le\_strongSolutionGrowthFloor\_le}, assembled from
\texttt{mainGrowthFloor\_le\_strongSolutionGrowthFloor} (the eigenpairs of
Theorem~\ref{thm:main-spectral} at the witnesses of Definition~\ref{def:main:witnesses})
and \texttt{reducedInductionOperator\_strongSolutionGrowthFloor\_le} (the ceiling, fed by
the a priori bound of Lemma~\ref{lem:dynamics:energy-bound}), both in
\texttt{FastDynamo/Dynamics/MainGrowthFloor.lean}.  So
$\gamma_0\leq\gamma_{\rm str}(A_{\eps,K_\eps})\leq M_u $ holds for every $0<\eps\leq\eps_0 $:
it is the counterpart of \eqref{eq:main-growth-conclusion} for $\gamma_{\rm str} $, the lower bound carrying the growth conclusion and the
upper bound showing that $\gamma_{\rm str} $ is a real quantity and not a vacuous
$+\infty $.  It is a statement about the initial-value problem on the domain; the
corresponding statement about $\gamma(u,\eps) $ is Theorem~\ref{thm:main-growth-rate}, and
it needs the whole generation chain of the next three sections, not this lemma.

\begin{proof}
  \leanok
  \uses{lem:dynamics:eigen-trajectory-law, lem:dynamics:energy-bound, thm:main-spectral, def:main:witnesses}
  The eigen-trajectory $t\mapsto\e^{t\lambda}B $ of
  Lemma~\ref{lem:dynamics:eigen-trajectory-law} is a forward strong solution with
  $\iota(B(0))=\iota(B)\neq0 $ and
  $\norm{\iota(\e^{t\lambda}B)}/\norm{\iota(B)}=\e^{t\re\lambda} $ for every $t $.  Taking
  logarithms and dividing by $t>0 $ makes the inner integrand identically $\re\lambda $
  along $t\to\infty $, so the $\limsup $ for this single solution is $\re\lambda $, and the
  supremum over all solutions is at least that.  For the ceiling, the hypothesis bounds
  the inner integrand of every admissible solution by $t^{-1}\log\e^{Mt}=M $ along
  $t\to\infty $ (the times where $\iota(B(t)) $ vanishes contribute $-\infty $), so every
  inner $\limsup $, and hence the supremum, is at most $M $.  The instance for
  $A_{\eps,K} $ is Lemma~\ref{lem:dynamics:energy-bound} with $M=M_u $, and the lower
  bound at the witnesses of Definition~\ref{def:main:witnesses} is the eigenpair clause of
  Theorem~\ref{thm:main-spectral} fed to the displayed eigenpair bound.
\end{proof}

\section{The Hille--Yosida generation theorem}
\label{sec:dynamics:hille-yosida}

Mathlib has no $\C_0 $-semigroup theory, so the generation theorem that item~(1) of
Lemma~\ref{lem:semigroup-spectral-growth} needs is built here from scratch, in
\texttt{FastDynamo/Dynamics/HilleYosida/}.  The paper generates an \emph{analytic}
semigroup by an infinitesimal perturbation of $\eps\Delta $; the formalisation instead uses
the dissipative route, whose only analytic input is a resolvent estimate on a real ray and
whose conclusion is a $\C_0 $-semigroup with an exponential norm bound.  Its generator is
identified in Lemma~\ref{lem:dynamics:hille-yosida-generator}, and analyticity is proved from a
sector condition in Theorem~\ref{thm:dynamics:analytic-generation}, on the same Yosida
approximants.

\begin{definition}[Hille--Yosida ray]
  \label{def:dynamics:hille-yosida-ray}
  \lean{FastDynamo.HilleYosidaRay, FastDynamo.HilleYosidaRay.ofHalfPlane,
    FastDynamo.HilleYosidaRay.norm_smul_stateResolvent_le,
    FastDynamo.HilleYosidaRay.tendsto_smul_stateResolvent}
  \leanok
  \uses{def:main:induction-operator}
  Let $A $ be an operator presented as a bounded map out of a Banach domain $D $ with
  continuous injection $\iota:D\to X $, and let $M\in\R $.  A \emph{Hille--Yosida ray of
  constant $M $} for $A $ is a set $S\subset\CC $ together with a two-sided inverse
  $(z-A)^{-1}:X\to D $ of the pencil $z\iota-A $ at every $z\in S $, such that $S $ contains
  every real $r>M $ and the state resolvent $R(r)=\iota\circ(r-A)^{-1} $ obeys
  \[
    \norm{R(r)}_{X\to X}\leq(r-M)^{-1}
    \qquad(r\in\R,\ r>M).
  \]
\end{definition}

In Lean this is the structure \texttt{FastDynamo.HilleYosidaRay A M}
(\texttt{FastDynamo/Dynamics/HilleYosida/ResolventRay.lean}), with fields
\texttt{spectralSet}, a \texttt{HasResolventOn} certificate, membership of the real ray and
the displayed bound.  The certificate records only the two inverse equations: no closed
graph, no completeness of $D $ and no continuity of $z\mapsto R(z) $ is assumed, and every
algebraic identity the construction needs (the resolvent identity, $AR(z)=zR(z)-\id $,
commutation of the resolvents) is derived from them.  \texttt{HilleYosidaRay.ofHalfPlane}
builds a ray from a resolvent on the half-plane $\{\re z>M\} $, which is how the two
instances below arise.  The two quantitative consequences the approximation scheme consumes
are \texttt{norm\_smul\_stateResolvent\_le} ($\norm{rR(r)}\leq r/(r-M) $, a bound uniform for
large $r $) and, when $\iota $ has dense range, the strong convergence
$rR(r)x\to x $ of \texttt{tendsto\_smul\_stateResolvent}.

\begin{theorem}[Hille--Yosida generation]
  \label{thm:dynamics:hille-yosida}
  \lean{FastDynamo.exists_stronglyContinuousSemigroup_of_hilleYosidaRay,
    FastDynamo.HilleYosidaRay.semigroup, FastDynamo.HilleYosidaRay.realization,
    FastDynamo.HilleYosidaRay.semigroup_spec,
    FastDynamo.HilleYosidaRay.norm_semigroup_le,
    FastDynamo.HilleYosidaRay.growthBound_le}
  \leanok
  \uses{def:dynamics:hille-yosida-ray, def:dynamics:semigroup-interface,
    def:dynamics:growth-bound}
  Let $X $ be a complex Banach space and let $A $ carry a Hille--Yosida ray of constant $M $
  whose domain injection $\iota:D\to X $ has dense range.  Then $A $ generates a strongly
  continuous semigroup $S $ on $X $: there is a realization of $S $ by $A $ in the sense of
  Definition~\ref{def:dynamics:semigroup-interface} --- every $d\in D $ evolves along a
  domain-valued forward strong solution whose embedding is the orbit $t\mapsto S(t)\iota(d) $
  --- and
  \[
    \norm{S(t)}_{X\to X}\leq\e^{Mt}
    \qquad(t\geq0),
    \qquad\text{hence}\qquad
    \omega_0(S)\leq M .
  \]
\end{theorem}

The existential form is
\texttt{exists\_stronglyContinuousSemigroup\_of\_hilleYosidaRay}; because the consumers need
the orbit and not only its existence, the same content is also stated about the named
objects \texttt{HilleYosidaRay.semigroup} and \texttt{HilleYosidaRay.realization}, as
\texttt{HilleYosidaRay.semigroup\_spec} (\texttt{FastDynamo/Dynamics/HilleYosida/Generation.lean}).
This is the analogue of the classical generation theorem for contraction-type semigroups
[Engel--Nagel, Theorem II.3.5] in this development's \texttt{DomainOperator} formulation.
Analyticity is not part of this theorem: it needs the additional sector condition of
Theorem~\ref{thm:dynamics:analytic-generation}.

\begin{proof}
  \leanok
  \uses{def:dynamics:hille-yosida-ray, def:dynamics:semigroup-interface,
    lem:dynamics:growth-bound-norm-bounds}
  The Yosida route.  The approximants $A_r:=r^2R(r)-r $
  (\texttt{HilleYosidaRay.yosida}) are bounded operators with
  $\norm{\e^{tA_r}}\leq\e^{Mrt/(r-M)} $, so their exponentials
  $\e^{tA_r} $ (\texttt{HilleYosidaRay.approx}) are uniformly bounded on compact time
  intervals; $A_r\iota(d)\to A d $ for every $d\in D $, and the two approximants commute,
  which turns the telescoping identity
  $\e^{tA_r}x-\e^{tA_s}x=\int_0^t\frac{\dd}{\dd\tau}\bigl(\e^{(t-\tau)A_s}\e^{\tau A_r}x\bigr)\dd\tau $
  into the Cauchy estimate of \texttt{cauchySeq\_approx}, first on the dense range of $\iota $
  and then, by the uniform bound, on all of $X $.  Completeness of $X $ gives the strong
  limit $S(t) $ (\texttt{HilleYosidaRay.semigroupOperator}); the semigroup law and
  $S(0)=\id $ pass to the limit, and strong continuity at $t=0 $ comes from the convergence
  being uniform on compact time intervals (\texttt{tendstoUniformlyOn\_approx}).  The
  exponential bound survives the limit, and the majorant half of
  Lemma~\ref{lem:dynamics:growth-bound-norm-bounds} turns it into $\omega_0(S)\leq M $.  For
  the realization, fix a real $\lambda_0>M $ and transport the ambient orbit of
  $(\lambda_0-A)d $ back to $D $ along $(\lambda_0-A)^{-1} $; because $S(t) $ commutes with
  $R(\lambda_0) $, the embedding of that domain-valued curve is $S(t)\iota(d) $ and the
  operator applied to it is $S(t)(Ad) $.  Its one-sided derivative at each $t\geq0 $ then
  comes from the integral identity
  $S(t)\iota(d)-\iota(d)=\int_0^tS(\tau)(Ad)\dd\tau $ by the second fundamental theorem of
  calculus, the endpoint $t=0 $ included.
\end{proof}

\begin{lemma}[The generated semigroup is the only one]
  \label{lem:dynamics:semigroup-unique}
  \lean{FastDynamo.StronglyContinuousSemigroup.operator_eq_of_realization,
    FastDynamo.HilleYosidaRay.semigroup_unique}
  \leanok
  \uses{def:dynamics:semigroup-interface, thm:dynamics:hille-yosida}
  Let $A $ have a densely embedded domain and unique forward strong solutions.  Then two
  strongly continuous semigroups admitting a realization by $A $ have the same evolution
  operator at every time.  In particular, if $A $ carries a Hille--Yosida ray, every
  strongly continuous semigroup realizing $A $ equals the one built in
  Theorem~\ref{thm:dynamics:hille-yosida}.
\end{lemma}

This is what earns the definite article in ``\emph{the} semigroup of $A $'' below, and it
removes any content from the choices made inside the construction: the semigroup does not
depend on which resolvent certificate the ray carries.  In Lean:
\texttt{StronglyContinuousSemigroup.operator\_eq\_of\_realization}
(\texttt{FastDynamo/Dynamics/StronglyContinuousSemigroup.lean}) and
\texttt{HilleYosidaRay.semigroup\_unique}
(\texttt{FastDynamo/Dynamics/HilleYosida/Generation.lean}).

\begin{proof}
  \leanok
  \uses{def:dynamics:semigroup-interface}
  On the image of $\iota $ the two realizations give two forward strong solutions with the
  same initial state, so forward uniqueness makes the two orbits agree there; both evolution
  operators are bounded and the image of $\iota $ is dense, so they agree on $X $.  As in
  Lemma~\ref{lem:dynamics:conditional-growth-bound}, forward uniqueness is a hypothesis of
  the statement rather than a lemma the proof calls: the concrete instance
  Lemma~\ref{lem:dynamics:strong-solution-uniqueness} lives in a Lean module that
  \emph{imports} this one, so it is not a \texttt{\textbackslash uses} of this proof.
\end{proof}

\begin{lemma}[Growth-bound sandwich for a Hille--Yosida generator]
  \label{lem:dynamics:hille-yosida-growth-sandwich}
  \lean{FastDynamo.re_le_growthBound_of_hilleYosidaRay_of_eigenpair}
  \leanok
  \uses{thm:dynamics:hille-yosida, def:dynamics:growth-bound}
  Let $A $ carry a Hille--Yosida ray of constant $M $, let its domain embed densely, and let
  its forward strong solutions be unique.  Then for every eigenpair $Ab=\lambda\iota(b) $
  with $\iota(b)\neq0 $,
  \[
    \re\lambda\ \leq\ \omega_0(S)\ \leq\ M ,
  \]
  $S $ being the semigroup of Theorem~\ref{thm:dynamics:hille-yosida}.
\end{lemma}

This is the shape in which the rest of the chapter uses the generation theorem: it converts
an eigenvalue of the operator into a two-sided pin on the growth bound of a
constructed semigroup.  It is
\texttt{re\_le\_growthBound\_of\_hilleYosidaRay\_of\_eigenpair}.

\begin{proof}
  \leanok
  \uses{lem:dynamics:conditional-growth-bound, thm:dynamics:hille-yosida}
  The upper bound is the exponential norm bound of Theorem~\ref{thm:dynamics:hille-yosida}.
  For the lower bound, the realization and forward uniqueness put the semigroup orbit of the
  eigenvector equal to the eigen-trajectory $t\mapsto\e^{t\lambda}b $, which is exactly the
  hypothesis discharged in Lemma~\ref{lem:dynamics:conditional-growth-bound}; its conclusion
  $\re\lambda\leq\omega_0(S) $ applies verbatim.
\end{proof}

\begin{lemma}[Invariant closed subspaces and the restricted semigroup]
  \label{lem:dynamics:invariant-restriction}
  \lean{FastDynamo.HilleYosidaRay.semigroup_mem_of_closed_resolvent_invariant,
    FastDynamo.HilleYosidaRay.restrictedSemigroup,
    FastDynamo.StronglyContinuousSemigroup.restrict,
    FastDynamo.StronglyContinuousSemigroup.growthBound_restrict_le,
    FastDynamo.StronglyContinuousSemigroup.congr,
    FastDynamo.StronglyContinuousSemigroup.growthBound_congr,
    FastDynamo.StronglyContinuousSemigroup.growthBound_mono_of_norm_le}
  \leanok
  \uses{thm:dynamics:hille-yosida, def:dynamics:growth-bound}
  Let $A $ carry a Hille--Yosida ray of constant $M $ and a densely embedded domain, and let
  $W\subset X $ be a closed subspace such that $R(r)W\subset W $ for every real $r>M $.  Then
  the generated semigroup leaves $W $ invariant, and its restriction $S|_W $ is a strongly
  continuous semigroup on $W $ with $\norm{S|_W(t)}\leq\e^{Mt} $ and
  $\omega_0(S|_W)\leq\omega_0(S) $.  Moreover, conjugating a strongly continuous semigroup by
  a surjective linear isometry changes no operator norm and hence no growth bound, and
  $\norm{S(t)}\leq\norm{T(t)} $ for all $t\geq0 $ implies $\omega_0(S)\leq\omega_0(T) $.
\end{lemma}

The invariance hypothesis is checked on the resolvent -- an explicit operator -- and
propagated to the semigroup along the approximation scheme,
$R(r)\rightsquigarrow A_r\rightsquigarrow\e^{tA_r}\rightsquigarrow S(t) $; closedness of $W $
is used exactly once, to pass the exponential series to the limit.  This is the only
available route, since the generated semigroup is a limit and not a formula.  Note the
direction: passing to an invariant subspace can strictly \emph{lower} a growth bound, so the
reverse inequality is false in general and is never claimed; a lower bound on
$\omega_0(S|_W) $ has to be proved on $W $ itself.  In Lean:
\texttt{FastDynamo/Dynamics/HilleYosida/Invariance.lean} for the invariance and
\texttt{HilleYosidaRay.restrictedSemigroup}, and
\texttt{FastDynamo/Dynamics/SemigroupRestriction.lean} for
\texttt{StronglyContinuousSemigroup.restrict}, \texttt{congr} and the two growth-bound
comparisons.

\begin{proof}
  \leanok
  \uses{thm:dynamics:hille-yosida, def:dynamics:growth-bound}
  A closed subspace invariant under a bounded operator $B $ is invariant under $\e^B $,
  because the exponential series converges in operator norm and every partial sum lies in
  $W $; the Yosida approximant $A_r=r^2R(r)-r $ preserves $W $ as soon as $R(r) $ does on the
  whole ray, hence so does $\e^{tA_r} $, and $S(t)x $ is the limit of $\e^{tA_r}x $.  The
  restricted family is the ambient one read through the subtype inclusion, so the three
  semigroup laws and strong continuity are inherited and $\norm{S|_W(t)}\leq\norm{S(t)} $.
  The comparison of growth bounds is monotonicity of the $\limsup $, the times with
  $\norm{S(t)}=0 $ contributing $-\infty $; for a surjective isometry both conjugations
  preserve norms, so the comparison holds in both directions and the growth bounds are equal.
\end{proof}

\section{The induction operator on one axial sector}

The Hille--Yosida ray of each axial sector comes from the dissipativity proved for the
reduced fixed-sector operator $A_{\eps,K} $ of \eqref{eq:main:torus-sector}.  What dissipativity of
$A_{\eps,K} $ alone gives is the lower bound $(\re z-M_u)\norm{\iota b}\leq\norm{(z-A)b} $,
that is, injectivity with closed range; surjectivity of the pencil is the Lumer--Phillips
half, and it comes from the same bound for the formal adjoint.

\begin{lemma}[The displayed formal adjoint of the reduced operator]
  \label{lem:dynamics:reduced-adjoint}
  \lean{FastDynamo.reducedInductionAdjointOperator,
    FastDynamo.reducedInductionStretchingTranspose,
    FastDynamo.reducedInductionOperator_formalAdjointPairing,
    FastDynamo.re_inner_reducedInductionAdjointAction_le,
    FastDynamo.reducedInductionAdjointOperator_forwardStrongSolutionsUnique,
    FastDynamo.reducedInductionOperator_isAdjointRealization}
  \leanok
  \uses{def:main:induction-operator, lem:dynamics:energy-bound}
  For a transverse divergence-free $u\in\W^{1,\infty} $, $\eps\geq0 $ and $K\in\Z $ put
  \[
    A^*_{\eps,K}b
    =\eps\bigl(\Delta_\perp-(2\pi K)^2\bigr)b
    +u_\perp\cdot\nabla_\perp b+2\pi iKu_3b
    +(\nabla_\perp u)^{\mathsf T}b,
    \qquad D(A^*_{\eps,K})=\Hs^2(\T^2;\CC^3),
  \]
  on the same domain and with the same $\Lp^2 $ embedding as $A_{\eps,K} $.  Then
  $\inner{A_{\eps,K}b}{c}=\inner{b}{A^*_{\eps,K}c} $ for all $b,c $ in the domain, the
  numerical range of $A^*_{\eps,K} $ lies in the same half-plane,
  \[
    \re\inner{b}{A^*_{\eps,K}b}\leq M_u\norm{b}_2^2,
  \]
  and $A^*_{\eps,K} $ has unique forward strong solutions.  If moreover $\eps>0 $, then
  $A^*_{\eps,K} $ is a maximal adjoint realization of $A_{\eps,K} $, and $A_{\eps,K} $ has a
  closed graph.
\end{lemma}

Two devices let this rest on facts proved elsewhere.  The transposed stretching matrix
is the pointwise adjoint $\star $ of the literal stretching matrix field of
\texttt{FastDynamo/Global/TorusLocalCoreMeasurability.lean}, so no new $\W^{1,\infty} $ datum is
introduced; and the three bilinear integration-by-parts identities (diffusion, transverse
transport, axial transport) are obtained from the \emph{quadratic} identities proved
in \texttt{FastDynamo/Global/ReducedInductionEnergyIdentity.lean} by complex polarization,
which over $\CC $ loses nothing.  In particular the divergence-freeness of $u $ enters
exactly where it enters the energy identity of Lemma~\ref{lem:dynamics:energy-bound}.  The
passage from the formal pairing to a \emph{maximal} adjoint realization, and with it the
closed graph, needs one far real reference point at which both $A_{\eps,K} $ and
$A^*_{\eps,K} $ are invertible; that is where $\eps>0 $ is used.  Module
\texttt{FastDynamo/Global/ReducedInductionAdjoint.lean}, with the reference point and the
maximality in \texttt{FastDynamo/Global/ReducedInductionHalfPlaneResolvent.lean}.

\begin{proof}
  \leanok
  \uses{lem:dynamics:energy-bound}
  A sesquilinear form over $\CC $ vanishing on the diagonal vanishes identically (test with
  $b+c $ and $b+ic $), so each quadratic identity
  $\im\inner{b}{\eps\Delta_\perp b}=0 $, $\re\inner{b}{u_\perp\cdot\nabla_\perp b}=0 $,
  $\re\inner{b}{2\pi iKu_3b}=0 $ upgrades to the corresponding bilinear pairing identity, and
  the stretching term contributes the pointwise adjoint relation
  $\inner{(b_\perp\cdot\nabla_\perp)u}{c}=\inner{b}{(\nabla_\perp u)^{\mathsf T}c} $
  integrated over $\T^2 $.  Summing the four gives the formal pairing.  The half-plane bound
  for $A^*_{\eps,K} $ is then the same computation as for $A_{\eps,K} $, the diffusion term
  being nonpositive and the transposed stretching matrix having the same $\Lp^\infty $ norm;
  uniqueness of forward strong solutions follows from
  Lemma~\ref{lem:dynamics:strong-solution-uniqueness} applied to $A^*_{\eps,K} $.  For the
  last claim, the exponential-ray Neumann construction of
  \texttt{FastDynamo/Operator/FixedSectorLowerOrderDecay.lean} inverts the diffusion pencil
  far out on the positive real axis and absorbs the lower-order part; it applies verbatim to
  $A^*_{\eps,K} $, whose diffusion pencil is the same and whose lower-order part decays
  relative to the diffusion inverse in the same way, the transposed stretching matrix being
  an $\Lp^\infty $ multiplier composed with the $\Hs^2\to\Lp^2 $ embedding.
\end{proof}

\begin{proposition}[The reduced operator has a resolvent on the right half-plane]
  \label{prop:dynamics:reduced-half-plane-resolvent}
  \lean{FastDynamo.reducedInductionRightHalfPlane,
    FastDynamo.reducedInductionOperator_pencil_bijective_of_lt_re,
    FastDynamo.reducedInductionRightHalfPlaneResolvent,
    FastDynamo.norm_stateResolvent_reducedInductionOperator_rightHalfPlane_le,
    FastDynamo.norm_stateResolvent_reducedInductionOperator_realRay_le}
  \leanok
  \uses{lem:dynamics:reduced-adjoint, lem:dynamics:energy-bound}
  Let $u $ be transverse divergence-free, $\eps>0 $ and $K\in\Z $.  Then for every
  $z\in\CC $ with $\re z>M_u $ the pencil $z-A_{\eps,K}:\Hs^2(\T^2;\CC^3)\to\Lp^2 $ is
  bijective, and the resolvents assemble into a single continuous family on the open
  half-plane $\{\re z>M_u\} $ with
  \[
    \norm{(z-A_{\eps,K})^{-1}}_{\Lp^2\to\Lp^2}\leq(\re z-M_u)^{-1} .
  \]
  In particular the bound holds at every real $r>M_u $.
\end{proposition}

This is the Lumer--Phillips input of Theorem~\ref{thm:dynamics:sector-semigroup}: the real-ray
specialization \texttt{norm\_stateResolvent\_reducedInductionOperator\_realRay\_le} is
exactly the hypothesis of Definition~\ref{def:dynamics:hille-yosida-ray}.  Module
\texttt{FastDynamo/Global/ReducedInductionHalfPlaneResolvent.lean}; the set
$\{\re z>M_u\} $ is \texttt{reducedInductionRightHalfPlane u}.  The coherent family is
assembled from the pointwise inverses by the repository's rational-propagation construction
\texttt{DomainOperator.hasResolventOnOfPointwiseTwoSidedInverses}, following the template
already used for the periodic background, so no continuity in $z $ is proved by hand.

\begin{proof}
  \leanok
  \uses{lem:dynamics:reduced-adjoint, lem:dynamics:energy-bound}
  The numerical-range bound of Lemma~\ref{lem:dynamics:energy-bound} places the numerical
  range of $A_{\eps,K} $ in $\{\re\nu\leq M_u\} $, which gives the ambient pencil lower
  bound $(\re z-M_u)\norm{\iota b}\leq\norm{(z-A_{\eps,K})b} $; hence the pencil is injective
  with closed range, and the same bound at $\bar z $ for the adjoint
  (Lemma~\ref{lem:dynamics:reduced-adjoint}) makes the range dense.  Closed plus dense is
  everything, so the pencil is bijective, and the open mapping theorem gives a bounded
  inverse; dividing the lower bound by $\norm{(z-A_{\eps,K})b} $ turns it into the displayed
  estimate.  The closed graph of $A_{\eps,K} $ and the maximality of the adjoint, both from
  the far real reference point of Lemma~\ref{lem:dynamics:reduced-adjoint}, are what make
  the range closed and identify the displayed formal adjoint with the Hilbert adjoint.
\end{proof}

\begin{theorem}[The Hille--Yosida ray of one axial sector]
  \label{thm:dynamics:sector-semigroup}
  \lean{FastDynamo.reducedInductionHilleYosidaRay,
    FastDynamo.reducedInductionHilleYosidaRay_spectralSet,
    FastDynamo.norm_stateResolvent_reducedInductionHilleYosidaRay_le}
  \leanok
  \uses{prop:dynamics:reduced-half-plane-resolvent, def:dynamics:hille-yosida-ray}
  Let $u $ be transverse divergence-free, $\eps>0 $ and $K\in\Z $.  Then $A_{\eps,K} $
  carries a Hille--Yosida ray of constant $M_u $ in the sense of
  Definition~\ref{def:dynamics:hille-yosida-ray}, whose spectral set is the whole open
  half-plane $\{\re z>M_u\} $ and whose state resolvent obeys
  $\norm{(r-A_{\eps,K})^{-1}}_{\Lp^2\to\Lp^2}\leq(r-M_u)^{-1} $ at every real $r>M_u $.
\end{theorem}

The ray is the input that the next section sums over the axial sectors: its resolvent
certificate is the summand of the all-frequency resolvent of
Proposition~\ref{prop:dynamics:product-resolvent}.  No semigroup is generated at the level
of a single sector; the semigroup of the development is that of the all-frequency operator
(Theorem~\ref{thm:dynamics:product-semigroup}).  The spectral set being the whole
half-plane is definitional: \texttt{HilleYosidaRay.ofHalfPlane} keeps the half-plane
certificate and only adds the ray estimate.  Module
\texttt{FastDynamo/Dynamics/ReducedInductionHilleYosidaRay.lean}.

\begin{proof}
  \leanok
  \uses{prop:dynamics:reduced-half-plane-resolvent, def:dynamics:hille-yosida-ray}
  Restrict the half-plane resolvent of
  Proposition~\ref{prop:dynamics:reduced-half-plane-resolvent} to the real ray by
  \texttt{HilleYosidaRay.ofHalfPlane}; the ray estimate is the half-plane estimate at real
  $z=r $.
\end{proof}

\section{From the axial sectors to the solenoidal semigroup}
\label{sec:dynamics:solenoidal}

The paper's semigroup acts on $\Lp^2_\sigma(\T^3) $.  Three steps separate it from the
fixed-sector operator $A_{\eps,K} $ and its Hille--Yosida ray
(Theorem~\ref{thm:dynamics:sector-semigroup}): summing the axial
sectors into the all-frequency operator $\mathcal A_\eps $ on $\Lp^2(\T^3;\CC^3) $,
restricting to the closed solenoidal subspace, and transporting to the literal three-torus.
The first step is the only one with analytic content, and it is done at the level of
resolvents rather than of semigroups: the sector resolvents sum to a resolvent as soon as
one elliptic estimate is uniform in the axial index, and the generation theorem is then
applied once, to the sum.

\begin{lemma}[The axial sectors decompose product $\Lp^2$ and diagonalise the operator]
  \label{lem:dynamics:axial-hilbert-sum}
  \lean{FastDynamo.isHilbertSum_axialLiftL2, FastDynamo.axialDecomposition,
    FastDynamo.axialSectorProjL2, FastDynamo.hasSum_axialLiftL2_axialSectorProjL2,
    FastDynamo.hasSum_norm_sq_axialSector, FastDynamo.hasSum_re_inner_axialSector,
    FastDynamo.axialSectorProjH2, FastDynamo.norm_axialSectorProjH2_le,
    FastDynamo.hasSum_axialLiftH2_axialSectorProjH2,
    FastDynamo.axialSectorProjL2_productInductionAction}
  \leanok
  \uses{def:main:induction-operator}
  The isometries $b\mapsto\e^{2\pi iKx_3}b $, $K\in\Z $, make $\Lp^2(\T^3;\CC^3) $ the
  Hilbert sum of copies of $\Lp^2(\T^2;\CC^3) $: they have pairwise orthogonal ranges, and
  the closed span of those ranges is everything.  Writing $P_K $ for the resulting
  orthogonal projection onto the $K $-th sector, read as a transverse amplitude,
  \[
    \norm{V}_2^2=\sum_{K\in\Z}\norm{P_KV}_2^2
    \qquad(V\in\Lp^2(\T^3;\CC^3)).
  \]
  The same decomposition lifts to the $\Hs^2 $ domain: there are projections
  $\Pi_K:\Hs^2(\T^3;\CC^3)\to\Hs^2(\T^2;\CC^3) $ of norm at most one, uniformly in $K $,
  such that every $B\in\Hs^2(\T^3;\CC^3) $ is the $\Hs^2 $-convergent sum
  $B=\sum_K\e^{2\pi iKx_3}\Pi_KB $.  The all-frequency operator is block diagonal for this
  decomposition,
  \[
    P_K\bigl(\mathcal A_\eps B\bigr)=A_{\eps,K}\bigl(\Pi_KB\bigr)
    \qquad(K\in\Z,\ B\in\Hs^2(\T^3;\CC^3)).
  \]
\end{lemma}

The ambient statement is Fourier analysis on the product torus together with Mathlib's
\texttt{IsHilbertSum} API: orthogonality is Fubini plus orthonormality of the circle
characters, and totality is completeness of the product Fourier system.  The domain-level
projection is not the adjoint of the $\Hs^2 $ lift -- that lift is only
$1+(2\pi K)^2 $-bounded -- but the ambient sector projection of the stored field paired with
the $K $-th row of the stored Fourier certificate, rescaled by a ratio bounded by one; this
is what makes $\norm{\Pi_K}\leq1 $ uniform in $K $.  Modules
\texttt{FastDynamo/Dynamics/Assembly/AxialHilbertSum.lean} and
\texttt{FastDynamo/Dynamics/Assembly/AxialSectorDomain.lean}.

\begin{proof}
  \leanok
  \uses{def:main:induction-operator}
  For orthogonality, $\inner{\e^{2\pi ikx_3}b}{\e^{2\pi iKx_3}c} $ splits by Fubini into the
  transverse inner product times $\int_{\T}\e^{2\pi i(K-k)x_3}\dd x_3 $, which is $1 $ when
  $k=K $ and $0 $ otherwise.  For totality, a field orthogonal to every sector has every
  product Fourier coefficient zero, by testing against a transverse scalar field placed in
  one vector coordinate.  Parseval is then Mathlib's, and the domain-level sum follows from
  the $\ell^2 $ expansion of the stored certificate together with closedness of the Fourier
  graph submodule.  Block diagonality holds on a single lifted sector by the exact sector
  identity of Definition~\ref{def:main:induction-operator}, and both sides are continuous in
  $B $, so the $\Hs^2 $-convergent sector expansion extends it to the whole domain.
\end{proof}

\begin{lemma}[The all-frequency operator is dissipative up to $M_u$]
  \label{lem:dynamics:product-dissipative}
  \lean{FastDynamo.re_inner_productInductionAction_le,
    FastDynamo.productInductionOperator_forwardStrongSolutionsUnique}
  \leanok
  \uses{lem:dynamics:axial-hilbert-sum, lem:dynamics:energy-bound}
  For every transverse divergence-free $u\in\W^{1,\infty} $, every $\eps\geq0 $ and every
  $B\in\Hs^2(\T^3;\CC^3) $,
  \[
    \re\inner{B}{\mathcal A_\eps B}\leq M_u\norm{B}_2^2 ,
  \]
  and consequently $\mathcal A_\eps $ has unique forward strong solutions, with no further
  hypothesis.
\end{lemma}

The energy competition is never redone at the product level: transport skewness and the sign
of the diffusion are consumed only through the reduced operator, sector by sector, which is
why the constant is the sector constant $M_u $ and is independent of $\eps $ and of the axial
index.  The statement is about the auxiliary operator $\mathcal A_\eps $; uniqueness for the
$\Hs^{-2}$ curl-form realisation in which clause~(16) of
Definition~\ref{def:main:physical-clauses} is stated is
Lemma~\ref{lem:dynamics:physical-uniqueness}, which reduces to this one.  Module
\texttt{FastDynamo/Dynamics/Assembly/AxialSectorDomain.lean}.

\begin{proof}
  \leanok
  \uses{lem:dynamics:axial-hilbert-sum, lem:dynamics:energy-bound,
    lem:dynamics:strong-solution-uniqueness}
  Both sides are sums over the axial sectors: the left-hand side by block diagonality and
  Parseval for the real part of the inner product, the right-hand side by Parseval for the
  norm.  Termwise the inequality is the fixed-sector bound of
  Lemma~\ref{lem:dynamics:energy-bound} with the same constant $M_u $, so it survives the
  summation.  Uniqueness is then
  Lemma~\ref{lem:dynamics:strong-solution-uniqueness} applied to $\mathcal A_\eps $.
\end{proof}

\begin{lemma}[An elliptic estimate uniform in the axial index]
  \label{lem:dynamics:uniform-sector-estimate}
  \lean{FastDynamo.productH2Weight_le_diffusion_bound_times_symbol,
    FastDynamo.norm_axialMultiplier_mul_inverseSymbol_le,
    FastDynamo.exists_uniform_sigma_norm_lowerOrder_comp_diffusionInverse_le_half,
    FastDynamo.norm_axialLiftH2_le_graph}
  \leanok
  \uses{def:main:induction-operator}
  Let $\eps>0 $.  There is a threshold $\sigma $, as large as one likes, at which the
  lower-order part of $A_{\eps,K} $ composed with the inverse of the diffusion pencil at
  $\sigma $ has norm at most $\tfrac12 $ \emph{simultaneously for every} $K\in\Z $; and there
  is $C>0 $ with
  \[
    \norm{\e^{2\pi iKx_3}b}_{\Hs^2(\T^3)}
    \leq C\bigl(\norm{b}_{\Lp^2(\T^2)}+\norm{A_{\eps,K}b}_{\Lp^2(\T^2)}\bigr)
    \qquad(K\in\Z,\ b\in\Hs^2(\T^2;\CC^3)).
  \]
\end{lemma}

This is the analytic crux of the assembly, and the only place where anything beyond algebra
is needed.  The fixed-sector chain of \texttt{FastDynamo/Operator/FixedSectorLowerOrderDecay.lean}
produces both a threshold and an $\Hs^2 $ bound that degrade as $\abs K\to\infty $, because
the axial transport term carries $\norm{2\pi K} $ and the sector lift is only
$1+(2\pi K)^2 $-bounded.  Two $K $-free scalar estimates remove both dependencies: the
product-weight elliptic gain
$\texttt{productH2Weight}\,(n,K)\leq(\sigma^{-1}+\eps^{-1})\,\texttt{symbol}_{\eps,\sigma,K}(n) $,
the three-dimensional analogue of the transverse gain already used to land the diffusion
inverse in $\Hs^2 $; and the arithmetic--geometric bound
$2\pi\abs K\cdot\texttt{inverseSymbol}_{\eps,\sigma,K}(n)\leq(2\sqrt{\eps\sigma})^{-1} $,
which absorbs the axial factor into the diffusion symbol.  The uniform threshold is a separate construction alongside the fixed-sector one, which
the fixed-sector statements do not depend on.
Module \texttt{FastDynamo/Dynamics/Assembly/UniformSectorResolvent.lean}.

\begin{proof}
  \leanok
  \uses{def:main:induction-operator}
  The lower-order part of $A_{\eps,K} $ is transverse transport plus axial transport plus
  stretching.  The transverse-transport and stretching bounds relative to the diffusion
  inverse are already $K $-free; the axial one is made $K $-free by the
  arithmetic--geometric estimate above, and all three tend to zero as $\sigma\to\infty $, so
  a single $\sigma $ works for every sector.  At such a $\sigma $ the Neumann series
  converges, which bounds $\norm{b}_{\Hs^2} $ by the diffusion-pencil norm and hence, through
  the product-weight gain applied to the lifted field, the $\Hs^2(\T^3) $ norm of
  $\e^{2\pi iKx_3}b $ by the graph norm of $A_{\eps,K} $ with a constant independent of $K $.
\end{proof}

\begin{proposition}[The all-frequency resolvent and its Hille--Yosida ray]
  \label{prop:dynamics:product-resolvent}
  \lean{FastDynamo.productResolventInv,
    FastDynamo.productInductionHalfPlaneResolvent,
    FastDynamo.norm_stateResolvent_productInduction_le,
    FastDynamo.productHilleYosidaRay}
  \leanok
  \uses{prop:dynamics:reduced-half-plane-resolvent, lem:dynamics:axial-hilbert-sum,
    lem:dynamics:uniform-sector-estimate, def:dynamics:hille-yosida-ray}
  Let $u $ be transverse divergence-free and $\eps>0 $.  On the half-plane $\{\re z>M_u\} $
  the operator
  \[
    R(z)x:=\sum_{K\in\Z}\e^{2\pi iKx_3}\,(z-A_{\eps,K})^{-1}(P_Kx)
  \]
  converges in $\Hs^2(\T^3;\CC^3) $, is a two-sided inverse of the pencil
  $z-\mathcal A_\eps $, depends continuously on $z $, and obeys
  $\norm{R(z)}_{\Lp^2\to\Lp^2}\leq(\re z-M_u)^{-1} $.  Hence $\mathcal A_\eps $ carries a
  Hille--Yosida ray of constant $M_u $.
\end{proposition}

The constant is not recomputed at the product level: it is the sector constant $M_u $, which
is independent of $\eps $ and of the axial index, and Parseval turns the sector estimates
into the same estimate for the sum.  Module
\texttt{FastDynamo/Dynamics/Assembly/ProductResolvent.lean}.

\begin{proof}
  \leanok
  \uses{prop:dynamics:reduced-half-plane-resolvent, lem:dynamics:axial-hilbert-sum,
    lem:dynamics:uniform-sector-estimate}
  Square summability of the summands is where the uniform estimate is used: by
  Lemma~\ref{lem:dynamics:uniform-sector-estimate} and the identity
  $A_{\eps,K}(z-A_{\eps,K})^{-1}=z(z-A_{\eps,K})^{-1}-\id $, the $\Hs^2 $ norm of the
  $K $-th summand is at most a constant depending only on $z $, $\eps $ and $M_u $ times
  $\norm{P_Kx} $, and $\sum_K\norm{P_Kx}^2=\norm{x}^2 $.  A family of sector lifts with square-summable graph
  norms sums in $\Hs^2(\T^3;\CC^3) $, the physical halves through the Hilbert-sum
  decomposition of Lemma~\ref{lem:dynamics:axial-hilbert-sum} and the stored certificates
  through their $\ell^2 $ expansion, the two limits being glued by closedness of the Fourier
  graph submodule.  Both inverse equations hold sector by sector, by block diagonality, and
  transport to the sum by uniqueness of limits.  The Parseval bound on the sum gives the
  estimate.  Continuity in $z $ is not proved by hand: the pointwise inverses are fed to the
  same rational-propagation construction as in
  Proposition~\ref{prop:dynamics:reduced-half-plane-resolvent}, whose stored inverse is
  continuous and, two-sided pencil inverses being unique, agrees with the sector sum.
  Restricting to the real ray gives the Hille--Yosida ray.
\end{proof}

\begin{theorem}[The all-frequency induction semigroup]
  \label{thm:dynamics:product-semigroup}
  \lean{FastDynamo.productInductionSemigroup, FastDynamo.productInductionRealization,
    FastDynamo.norm_productInductionSemigroup_le,
    FastDynamo.growthBound_productInductionSemigroup_le,
    FastDynamo.re_le_growthBound_productInductionSemigroup_of_eigenpair,
    FastDynamo.productInductionSemigroup_unique}
  \leanok
  \uses{prop:dynamics:product-resolvent, thm:dynamics:hille-yosida,
    lem:dynamics:product-dissipative, lem:dynamics:hille-yosida-growth-sandwich,
    lem:dynamics:semigroup-unique}
  Let $u $ be transverse divergence-free and $\eps>0 $.  Then the auxiliary operator
  $\mathcal A_\eps $ with domain $\Hs^2(\T^3;\CC^3) $ generates a strongly continuous
  semigroup on $\Lp^2(\T^3;\CC^3) $, realizing it in the sense of
  Definition~\ref{def:dynamics:semigroup-interface}, with
  $\norm{\e^{t\mathcal A_\eps}}\leq\e^{M_ut} $, hence
  $\omega_0\leq M_u $, and with $\re\lambda\leq\omega_0 $ for every eigenpair of
  $\mathcal A_\eps $.  It is the only strongly continuous semigroup realizing
  $\mathcal A_\eps $ in this sense.
\end{theorem}

This is the generation step of item~(1) of Lemma~\ref{lem:semigroup-spectral-growth}, for the
\emph{auxiliary} operator on all of $\Lp^2(\T^3;\CC^3)$.  Uniqueness is
Lemma~\ref{lem:dynamics:semigroup-unique} applied to the realization
(\texttt{productInductionSemigroup\_unique}), so the choice of resolvent family inside the
construction is immaterial.  The paper's operator is the curl form on $\Lp^2_\sigma$; the two
ambient operators are different on $\Lp^2(\T^3;\CC^3)$ (they differ by $-u\,\Div B $, which
vanishes only on solenoidal fields), and their actions agree on divergence-free $\Hs^2 $
fields (\texttt{productInductionAction\_toHMinus2\_eq\_physical},
\texttt{FastDynamo/Operator/Compatibility.lean}).  The generator of this semigroup is
$\mathcal A_\eps$ with domain $\Hs^2$ (\texttt{generator\_productInductionSemigroup},
Lemma~\ref{lem:dynamics:hille-yosida-generator}), and the generator of its restriction to
$\Lp^2_\sigma$ is the part of $\mathcal A_\eps$ there, which is the paper's $L_{\eps,u}$
(Theorem~\ref{thm:dynamics:solenoidal-generator}).  Module
\texttt{FastDynamo/Dynamics/Assembly/ProductSemigroup.lean}.

\begin{proof}
  \leanok
  \uses{prop:dynamics:product-resolvent, thm:dynamics:hille-yosida,
    lem:dynamics:product-dissipative, lem:dynamics:hille-yosida-growth-sandwich}
  The Hille--Yosida ray is Proposition~\ref{prop:dynamics:product-resolvent}, the
  $\Hs^2\to\Lp^2 $ inclusion has dense range
  (\texttt{FastDynamo/Analysis/ProductH2DenseEmbedding.lean}), and forward strong solutions
  are unique by Lemma~\ref{lem:dynamics:product-dissipative}.  Apply
  Theorem~\ref{thm:dynamics:hille-yosida} and
  Lemma~\ref{lem:dynamics:hille-yosida-growth-sandwich}.
\end{proof}

\begin{lemma}[The solenoidal subspace is closed and resolvent-invariant]
  \label{lem:dynamics:solenoidal-subspace}
  \lean{FastDynamo.productSolenoidalSubspace, FastDynamo.mem_productSolenoidalSubspace_iff,
    FastDynamo.isClosed_productSolenoidalSubspace,
    FastDynamo.stateResolvent_mem_productSolenoidalSubspace_of_re_pos,
    FastDynamo.productH2Embedding_mem_productSolenoidalSubspace_of_isEigenpair}
  \leanok
  \uses{def:main:induction-operator, lem:positive-eigenfunction-solenoidal}
  The set $\Lp^2_\sigma(\T^3)=\{V\in\Lp^2(\T^3;\CC^3):\Div V=0\} $ is a closed subspace of
  $\Lp^2(\T^3;\CC^3) $, hence a Hilbert space in its own right.  It is invariant under every
  state resolvent of $\mathcal A_\eps $ at a spectral parameter with positive real part: if
  $\eps\geq0 $, $u $ is transverse divergence-free, $(z-\mathcal A_\eps)^{-1} $ exists,
  $\re z>0 $ and $\Div x=0 $, then $\Div\bigl((z-\mathcal A_\eps)^{-1}x\bigr)=0 $.  Moreover
  every eigenfunction of $\mathcal A_\eps $ with eigenvalue of positive real part lies in
  $\Lp^2_\sigma(\T^3) $.
\end{lemma}

Both halves are cheap here, and neither is a new analytic input.  Solenoidality is a family
of scalar linear equations in the Fourier coefficients -- for each transverse frequency $n $
and axial frequency $k $ the combination
$\sum_{j<2}2\pi in_j\widehat V_j(n,k)+2\pi ik\widehat V_3(n,k) $ is a continuous linear
functional -- so $\Lp^2_\sigma $ is an intersection of kernels, exactly as the Fourier graph
domain is.  Invariance is the resolvent form of
Lemma~\ref{lem:positive-eigenfunction-solenoidal}: taking the weak divergence of the
resolvent equation turns it into a scalar advection--diffusion equation for $d=\Div B $ whose
forcing is the weak divergence of the solenoidal datum and therefore vanishes, and the same
energy endpoint forces $d=0 $.  The statement is quantified over an \emph{arbitrary}
resolvent certificate, so it does not depend on the construction of
Proposition~\ref{prop:dynamics:product-resolvent}.  Module
\texttt{FastDynamo/Dynamics/Assembly/SolenoidalSubspace.lean}.

\begin{proof}
  \leanok
  \uses{lem:positive-eigenfunction-solenoidal}
  A submodule cut out by a family of continuous linear functionals is closed, and a closed
  subspace of a Hilbert space is complete.  For invariance, write $B=(z-\mathcal A_\eps)^{-1}x $
  and apply the weak divergence to $zB-\mathcal A_\eps B=x $: incompressibility of $u $ makes
  the transport and stretching terms combine into $-(u\cdot\nabla)\Div B $, so
  $d=\Div B $ satisfies $(z-\eps\Delta+u\cdot\nabla)d=\Div x=0 $ weakly, and the energy
  endpoint of Lemma~\ref{lem:positive-eigenfunction-solenoidal} with $\re z>0 $ gives $d=0 $.
  The eigenfunction statement is the same argument with $x=0 $.
\end{proof}

\begin{theorem}[The induction semigroup on the solenoidal space]
  \label{thm:dynamics:solenoidal-semigroup}
  \lean{FastDynamo.productInductionSemigroup_mem_productSolenoidalSubspace,
    FastDynamo.solenoidalInductionSemigroup,
    FastDynamo.growthBound_solenoidalInductionSemigroup_le,
    FastDynamo.re_le_growthBound_solenoidalInductionSemigroup_of_eigenpair,
    FastDynamo.literalSolenoidalSubspace, FastDynamo.isClosed_literalSolenoidalSubspace,
    FastDynamo.exists_ne_zero_literalSolenoidalSubspace,
    FastDynamo.solenoidalToLiteral, FastDynamo.literalInductionSemigroup,
    FastDynamo.norm_literalInductionSemigroup_le,
    FastDynamo.growthBound_literalInductionSemigroup,
    FastDynamo.productInductionRealization_domainOrbit_mem_productSolenoidalSubspace,
    FastDynamo.productL2ToHMinus2_domainOrbit_eq_physical}
  \leanok
  \uses{thm:dynamics:product-semigroup, lem:dynamics:solenoidal-subspace,
    lem:dynamics:invariant-restriction, lem:main:literal-torus}
  Let $u $ be transverse divergence-free and $\eps>0 $.  Then $\e^{t\mathcal A_\eps} $ leaves
  $\Lp^2_\sigma(\T^3) $ invariant, and its restriction
  $T(t):=\e^{t\mathcal A_\eps}|_{\Lp^2_\sigma} $ is a strongly continuous semigroup on
  $\Lp^2_\sigma(\T^3) $ with $\norm{T(t)}\leq\e^{M_ut} $ and $\omega_0(T)\leq M_u $.  For
  every eigenpair $\mathcal A_\eps V=\lambda V $ whose eigenfunction lies in
  $\Lp^2_\sigma(\T^3) $ one has $\re\lambda\leq\omega_0(T) $.  All of this transports, with
  no change of any operator norm and hence no change of the growth bound, to the literal
  three-torus $\T^3=(\R/\Z)^3 $, where the solenoidal subspace is again closed.  Moreover
  the orbit of every solenoidal $\Hs^2 $ datum is a solenoidal $\Hs^2 $ field at every
  $t\geq0 $ and solves the paper's curl-form induction equation
  $\partial_tB=\operatorname{curl}(u\times B)+\eps\Delta B $ in $\Hs^{-2} $ there.
\end{theorem}

Three points of the passage to the subspace matter.

First, invariance is proved on the resolvent, at every real $r>M_u\geq0 $, and propagated
through the approximation scheme by Lemma~\ref{lem:dynamics:invariant-restriction}; there is
no other route, the generated semigroup being a limit and not a formula.  Second, the
spectral lower bound has to be \emph{re-proved} on the subspace: restriction can strictly
lower a growth bound, so the ambient bound of
Theorem~\ref{thm:dynamics:product-semigroup} is no evidence about the restricted one.  It is
re-proved by evaluating the restricted operator on an eigenvector that lies in
$\Lp^2_\sigma $ -- the subspace carries the ambient norm, so nothing is lost -- which is why
Lemma~\ref{lem:dynamics:solenoidal-subspace} is needed twice.  That same eigenvector shows
$\Lp^2_\sigma(\T^3)\neq0$ (\texttt{exists\_ne\_zero\_literalSolenoidalSubspace}): it is
non-zero by the definition of an eigenpair and solenoidal by the auxiliary solenoidality
theorem.  The statement is recorded because every clause of the ideal half of
Theorem~\ref{thm:main} is quantified over $B\in\Lp^2_\sigma(\T^3) $ and would hold verbatim
on the zero space; see Corollary~\ref{cor:ideal-growth}.

Third, \emph{$T(t)$ is the paper's $\e^{tL_{\eps,u}}$}: its generator is $L_{\eps,u}$ with
domain $\Hs^2\cap\Lp^2_\sigma$, and a $\C_0$-semigroup is determined by its generator
(Theorem~\ref{thm:dynamics:solenoidal-generator}).  What remains a deviation is the route: the
paper generates on $\Lp^2_\sigma$ directly, by perturbation of $\eps\Delta$, whereas this node
restricts the semigroup of the ambient auxiliary operator, and the generator is identified
afterwards.  The orbit of every solenoidal $\Hs^2 $ datum stays solenoidal and satisfies the
curl-form equation in $\Hs^{-2} $ at every $t\geq0 $
(\texttt{productL2ToHMinus2\_domainOrbit\_eq\_physical}), and it is the only solution of that
equation with that datum (Lemma~\ref{lem:dynamics:physical-uniqueness}).  So
$\gamma(u,\eps)$ (Definition~\ref{def:main:growth-rate}) is the growth bound of $T $.

Modules \texttt{FastDynamo/Dynamics/Assembly/ProductSemigroup.lean} and
\texttt{FastDynamo/Dynamics/Assembly/MainGrowthRate.lean}; the isometric identification of
the product model with the literal three-torus is the one of
Lemma~\ref{lem:main:literal-torus}.

\begin{proof}
  \leanok
  \uses{thm:dynamics:product-semigroup, lem:dynamics:solenoidal-subspace,
    lem:dynamics:invariant-restriction}
  The state resolvent of the product Hille--Yosida ray at a real $r>M_u $ preserves
  $\Lp^2_\sigma $ by Lemma~\ref{lem:dynamics:solenoidal-subspace}, the positivity of $r $
  being free since $M_u\geq0 $; the subspace is closed, so
  Lemma~\ref{lem:dynamics:invariant-restriction} propagates the invariance to the semigroup
  and produces the restriction with the same exponential bound.  For the lower bound,
  evaluate the restricted operator at time $t $ on an eigenvector $V\in\Lp^2_\sigma $: its
  orbit is $\e^{t\lambda}V $ by Theorem~\ref{thm:dynamics:product-semigroup}, so
  $\e^{t\re\lambda}\norm{V}\leq\norm{T(t)}\norm{V} $ and the minorant half of
  Lemma~\ref{lem:dynamics:growth-bound-norm-bounds} applies.  The transport is conjugation by
  a surjective linear isometry, which preserves operator norms exactly, so the growth bound
  is unchanged (Lemma~\ref{lem:dynamics:invariant-restriction}).
\end{proof}

\section{The generator}
\label{sec:dynamics:generator}

The semigroup of the previous section is built from an operator; this section proves that the
operator is its generator, which identifies the semigroup on $\Lp^2_\sigma$ with the paper's
$\e^{tL_{\eps,u}}$.

\begin{definition}[Generator, resolvent set and spectrum]
  \label{def:dynamics:generator}
  \lean{FastDynamo.StronglyContinuousSemigroup.diffQuotient,
    FastDynamo.StronglyContinuousSemigroup.generator,
    FastDynamo.StronglyContinuousSemigroup.mem_generator_domain_iff,
    FastDynamo.LinearPMap.IsResolvent, FastDynamo.LinearPMap.resolventSet,
    FastDynamo.LinearPMap.spectrum, FastDynamo.DomainOperator.toLinearPMap}
  \leanok
  \uses{def:dynamics:semigroup-interface}
  The \emph{generator} of a strongly continuous semigroup $S$ on $X$ is the operator
  \[
    Gx:=\lim_{h\downarrow0}\frac{S(h)x-x}{h},
    \qquad
    D(G)=\{x\in X:\text{the limit exists}\}.
  \]
  For a linear operator $T$ with domain $D(T)\subset X$, the \emph{resolvent set}
  $\rho(T)$ is the set of $z\in\CC$ at which $z-T:D(T)\to X$ has a bounded two-sided inverse
  $X\to X$ with range in $D(T)$, and the \emph{spectrum} is $\sigma(T)=\CC\setminus\rho(T)$.
  An operator presented by a bounded map out of a graph domain with an injective embedding
  (\texttt{DomainOperator}) is read as an operator on the range of the embedding.
\end{definition}

The definitions are the standard ones, which the paper uses without stating them.  In Lean the
generator is a \texttt{LinearPMap} (\texttt{FastDynamo/Dynamics/SemigroupGenerator.lean}), and
\texttt{LinearPMap.IsResolvent}, \texttt{resolventSet}, \texttt{spectrum} and
\texttt{DomainOperator.toLinearPMap} are in \texttt{FastDynamo/Operator/UnboundedResolvent.lean}.
The resolvent set asks for a bounded inverse explicitly; no closedness of $T$ is assumed.

\begin{lemma}[Properties of the generator]
  \label{lem:dynamics:generator-basic}
  \lean{FastDynamo.StronglyContinuousSemigroup.generator_eq_iff,
    FastDynamo.StronglyContinuousSemigroup.apply_mem_generator_domain,
    FastDynamo.StronglyContinuousSemigroup.hasDerivWithinAt_apply,
    FastDynamo.StronglyContinuousSemigroup.dense_generator_domain,
    FastDynamo.StronglyContinuousSemigroup.apply_eq_cexp_smul_of_generator_eq_smul,
    FastDynamo.StronglyContinuousSemigroup.ext_of_generator_eq,
    FastDynamo.StronglyContinuousSemigroup.mem_generator_restrict_domain_iff,
    FastDynamo.StronglyContinuousSemigroup.generator_restrict_apply,
    FastDynamo.LinearPMap.mem_spectrum_of_eigenvector,
    FastDynamo.LinearPMap.exists_eigenvector_of_mem_spectrum,
    FastDynamo.LinearPMap.IsResolvent.isCompactOperator_of_isResolvent,
    FastDynamo.DomainOperator.isResolvent_stateResolvent}
  \leanok
  \uses{def:dynamics:generator, lem:dynamics:invariant-restriction}
  Let $S$ be a strongly continuous semigroup on a complex Banach space $X$ with generator $G$.
  Then $D(G)$ is dense and invariant under $S(t)$, $GS(t)x=S(t)Gx$ on $D(G)$, and
  $t\mapsto S(t)x$ is differentiable on $[0,\infty)$ with derivative $GS(t)x$ for $x\in D(G)$;
  an eigenvector $Gx=\lambda x$ has $S(t)x=\e^{t\lambda}x$.  Two strongly continuous semigroups
  with the same generator are equal.  If $W\subset X$ is a closed subspace invariant under
  $S$, the generator of $S|_W$ is the part of $G$ in $W$: $D(G|_W)=D(G)\cap W$ with the same
  values.  Every eigenvalue of an operator lies in its spectrum; if the resolvent at one point
  is compact, it is compact at every point of the resolvent set, and then every point of the
  spectrum is an eigenvalue.
\end{lemma}

The semigroup facts are [Engel--Nagel, Lemma~II.1.3, Theorem~II.1.4, Proposition~II.2.3]; the
last two sentences are the resolvent identity and the Fredholm alternative for compact
operators.  The paper uses all of them without stating them.

\begin{proof}
  \leanok
  \uses{def:dynamics:generator}
  The derivative of $t\mapsto S(t)x$ is computed from the difference quotient on both sides;
  the integrals $\int_0^tS(s)x\dd s$ lie in $D(G)$ and approximate $x$, which gives density.
  For uniqueness, if $S$ and $T$ have the same generator then
  $s\mapsto T(t-s)S(s)x$ has derivative zero for $x\in D(G)$, so $S(t)x=T(t)x$ on a dense set
  and hence everywhere.  The restriction statement compares the difference quotients in $W$
  and in $X$.  For the spectrum: the resolvent identity
  $R(w)=R(z)+(z-w)R(w)R(z)$ transfers compactness, and for $z\in\sigma(T)$ and a compact
  resolvent $R(r)$ the operator $z-T$ is $(\id-(r-z)R(r))(r-T)$, so non-invertibility of
  $z-T$ forces a non-trivial kernel of $\id-(r-z)R(r)$ by the Fredholm alternative.
\end{proof}

\begin{lemma}[The generator of the Hille--Yosida semigroup]
  \label{lem:dynamics:hille-yosida-generator}
  \lean{FastDynamo.HilleYosidaRay.generator_semigroup,
    FastDynamo.HilleYosidaRay.tendsto_diffQuotient_semigroup_embedding,
    FastDynamo.generator_productInductionSemigroup}
  \leanok
  \uses{thm:dynamics:hille-yosida, def:dynamics:generator, lem:dynamics:generator-basic,
    thm:dynamics:product-semigroup}
  Let $A$ carry a Hille--Yosida ray whose domain embeds densely in the Banach space $X$, and let
  $S$ be the semigroup of Theorem~\ref{thm:dynamics:hille-yosida}.  Then the generator of $S$
  is $A$, with domain the range of the embedding.  In particular the generator of the
  all-frequency semigroup of Theorem~\ref{thm:dynamics:product-semigroup} is $\mathcal A_\eps$
  with domain $\Hs^2(\T^3;\CC^3)$.
\end{lemma}

This is [Engel--Nagel, Theorem~II.3.5] in the generator direction.  Modules
\texttt{FastDynamo/Dynamics/HilleYosida/GeneratorIdentification.lean} and
\texttt{FastDynamo/Dynamics/Assembly/SolenoidalGenerator.lean}.

\begin{proof}
  \leanok
  \uses{thm:dynamics:hille-yosida, lem:dynamics:generator-basic}
  The range of the embedding lies in $D(G)$ with $G=A$ there: the realization of
  Theorem~\ref{thm:dynamics:hille-yosida} differentiates the orbit at $t=0$.  Conversely, let
  $x\in D(G)$ and fix a real $r$ beyond the ray constant and with $\norm{S(1)}<\e^r$.  Put
  $w=R(r)(rx-Gx)-x$; then $w\in D(G)$ and $Gw=rw$, so $S(1)w=\e^rw$ by
  Lemma~\ref{lem:dynamics:generator-basic}, and $\norm{S(1)}<\e^r$ forces $w=0$.  Hence $x$
  is in the range of the embedding.
\end{proof}

\begin{theorem}[The induction semigroup is generated by $L_{\eps,u}$]
  \label{thm:dynamics:solenoidal-generator}
  \lean{FastDynamo.solenoidalH2, FastDynamo.productInductionAction_mem_productSolenoidalSubspace,
    FastDynamo.solenoidalInductionOperator,
    FastDynamo.generator_solenoidalInductionSemigroup,
    FastDynamo.mem_generator_solenoidalInductionSemigroup_domain_iff,
    FastDynamo.coe_generator_solenoidalInductionSemigroup,
    FastDynamo.productL2ToHMinus2_generator_solenoidalInductionSemigroup,
    FastDynamo.eq_solenoidalInductionSemigroup_of_generator_eq,
    FastDynamo.gamma_eq_growthBound_of_generator_eq,
    FastDynamo.exists_isResolvent_generator_solenoidalInductionSemigroup_eq}
  \leanok
  \uses{thm:dynamics:solenoidal-semigroup, lem:dynamics:hille-yosida-generator,
    lem:dynamics:generator-basic, lem:finish:physical-compatibility}
  Let $u$ be transverse divergence-free and $\eps>0$.  For every divergence-free
  $B\in\Hs^2$ the field $\mathcal A_\eps B$ lies in $\Lp^2_\sigma$, and it is the curl form
  $L_{\eps,u}B=\curl(u\times B)+\eps\Delta B$.  The generator of the semigroup $T$ of
  Theorem~\ref{thm:dynamics:solenoidal-semigroup} is $L_{\eps,u}$ with domain
  $\Hs^2\cap\Lp^2_\sigma$; every strongly continuous semigroup on $\Lp^2_\sigma$ with that
  generator is $T$, and its growth bound is $\gamma(u,\eps)$.
\end{theorem}

This is the generation statement of item~(1) of Lemma~\ref{lem:semigroup-spectral-growth},
without analyticity.  The Lean operator \texttt{solenoidalInductionOperator} has domain the
divergence-free $\Hs^2$ fields (\texttt{solenoidalH2}) and values the auxiliary action
$\mathcal A_\eps B$; the identification with the curl form is
\texttt{productL2ToHMinus2\_generator\_solenoidalInductionSemigroup}, an identity in the Fourier
model of $\Hs^{-2}$, into which $\Lp^2$ embeds injectively.  So the curl form of the paper's
\texttt{eq:physical-generator}, computed distributionally, is the $\Lp^2$ field
$\mathcal A_\eps B$: an equivalent formulation.  The route differs from the paper's: the
semigroup is generated on all of $\Lp^2$ by Hille--Yosida, its generator identified
(Lemma~\ref{lem:dynamics:hille-yosida-generator}), and restricted to $\Lp^2_\sigma$
(Lemma~\ref{lem:dynamics:generator-basic}, [Engel--Nagel, Proposition~II.2.3]), not generated
on $\Lp^2_\sigma$ by perturbation of $\eps\Delta$.  Module
\texttt{FastDynamo/Dynamics/Assembly/SolenoidalGenerator.lean}.

\begin{proof}
  \leanok
  \uses{lem:dynamics:hille-yosida-generator, lem:dynamics:generator-basic,
    lem:finish:physical-compatibility}
  On a divergence-free $\Hs^2$ field the lower-order part of $\mathcal A_\eps B$ is the weak
  curl $\curl(u\times B)$, which is solenoidal, and $\Delta B$ is solenoidal coefficientwise;
  the two actions agree in $\Hs^{-2}$ by Lemma~\ref{lem:finish:physical-compatibility}.  The
  generator of the restriction is the part in $\Lp^2_\sigma$ of the ambient generator
  $\mathcal A_\eps$ (Lemmas~\ref{lem:dynamics:hille-yosida-generator}
  and~\ref{lem:dynamics:generator-basic}), whose domain is
  $\Hs^2\cap\Lp^2_\sigma$.  Uniqueness is uniqueness of the semigroup with a given generator.
\end{proof}

\section{Analyticity, compactness and spectral mapping}
\label{sec:dynamics:analytic}

\begin{definition}[Analytic semigroup]
  \label{def:dynamics:analytic-semigroup}
  \lean{FastDynamo.openSector, FastDynamo.AnalyticExtension, FastDynamo.IsAnalyticSemigroup,
    FastDynamo.isOpen_openSector, FastDynamo.isPreconnected_openSector}
  \leanok
  \uses{def:dynamics:semigroup-interface}
  For $0<\delta\leq\pi/2$ let $\Sigma_\delta=\{z\in\CC\setminus\{0\}:\abs{\arg z}<\delta\}$.  A
  strongly continuous semigroup $S$ is \emph{analytic of angle $\delta$} if there is a map
  $T:\Sigma_\delta\cup[0,\infty)\to\mathcal L(X)$ that agrees with $S$ on $[0,\infty)$, obeys
  $T(z+w)=T(z)T(w)$ on $\Sigma_\delta$, is holomorphic on $\Sigma_\delta$ in the operator norm,
  and converges strongly to the identity as $z\to0$ within every smaller sector
  $\Sigma_{\delta'}$.
\end{definition}

This is [Engel--Nagel, Definition~II.4.5], the notion of the paper's
Lemma~\ref{lem:semigroup-spectral-growth}, with the angle explicit.  Module
\texttt{FastDynamo/Dynamics/AnalyticSemigroup.lean}.

\begin{lemma}[The sector condition]
  \label{lem:dynamics:sectorial}
  \lean{FastDynamo.productInductionTransport, FastDynamo.productInductionAction_eq_transport,
    FastDynamo.re_inner_productInductionAction_le_sub,
    FastDynamo.exists_norm_productInductionTransport_le,
    FastDynamo.re_inner_productInductionTransport_skew,
    FastDynamo.exists_norm_productInductionStretching_le,
    FastDynamo.exists_productInduction_sectorial}
  \leanok
  \uses{def:main:induction-operator, lem:dynamics:energy-identity}
  Let $u$ be transverse divergence-free and $\eps>0$.  There is $\omega\in\R$ with
  \[
    \abs{\im\inner{B}{\mathcal A_\eps B}}\leq\omega\norm B_2^2-\re\inner{B}{\mathcal A_\eps B}
    \qquad(B\in\Hs^2(\T^3;\CC^3)).
  \]
\end{lemma}

The paper states no such estimate; it is the numerical-range condition from which the Lean
proof derives analyticity, in place of the paper's relative bound for $\curl(u\times B)$.  The
constant is $\omega=M_u+d+c^2/(4\eps)$, with $c$ a bound of the transport term by the $\Hs^1$
seminorm and $d$ a bound of the stretching term in $\Lp^2$.  Module
\texttt{FastDynamo/Dynamics/Assembly/ProductInductionEnergy.lean}.

\begin{proof}
  \leanok
  \uses{lem:dynamics:energy-identity}
  The diffusion contributes the real number $-\eps\norm{\nabla B}_2^2$, transport is skew and
  contributes only to the imaginary part, bounded by $c\norm{\nabla B}_2\norm B_2$, and
  stretching contributes at most $d\norm B_2^2$ to either part.  Young's inequality absorbs
  $c\norm{\nabla B}_2\norm B_2$ into $\eps\norm{\nabla B}_2^2+\tfrac{c^2}{4\eps}\norm B_2^2$,
  and the energy identity bounds $\re\inner{B}{\mathcal A_\eps B}$ by
  $M_u\norm B_2^2-\eps\norm{\nabla B}_2^2$.
\end{proof}

\begin{theorem}[Analytic generation from the sector condition]
  \label{thm:dynamics:analytic-generation}
  \lean{FastDynamo.HilleYosidaRay.sectorial_yosida,
    FastDynamo.HilleYosidaRay.sectorExtension,
    FastDynamo.HilleYosidaRay.isAnalyticSemigroup_of_sectorial,
    FastDynamo.IsAnalyticSemigroup.restrict,
    FastDynamo.differentiableOn_of_forall_apply}
  \leanok
  \uses{thm:dynamics:hille-yosida, def:dynamics:analytic-semigroup,
    lem:dynamics:invariant-restriction}
  Let $X$ be a complex Hilbert space and let $A$ carry a Hille--Yosida ray with densely
  embedded domain.  If there is $\omega\in\R$ with
  $\abs{\im\inner{\iota d}{Ad}}\leq\omega\norm{\iota d}^2-\re\inner{\iota d}{Ad}$ for every
  $d$ in the domain, then the semigroup of Theorem~\ref{thm:dynamics:hille-yosida} is analytic
  of angle $\pi/4$.  The restriction of an analytic semigroup of angle $\delta$ to a closed
  invariant subspace is analytic of angle $\delta$.  A family of bounded operators on an open
  set of $\CC$ that is locally bounded and holomorphic at every vector is holomorphic in the
  operator norm.
\end{theorem}

The paper obtains analyticity from [Engel--Nagel, Theorem~III.2.10] (relative boundedness with
respect to $\eps\Delta$); this theorem replaces that step, and exhibits the angle.  Modules
\texttt{FastDynamo/Dynamics/HilleYosida/AnalyticGeneration.lean},
\texttt{FastDynamo/Dynamics/AnalyticSemigroup.lean} and
\texttt{FastDynamo/Operator/StrongHolomorphy.lean}.

\begin{proof}
  \leanok
  \uses{thm:dynamics:hille-yosida, lem:dynamics:invariant-restriction}
  The sector condition passes to the Yosida approximants $A_s=s^2R(s)-s$ of the proof of
  Theorem~\ref{thm:dynamics:hille-yosida}, with $\abs\omega+1$ in place of $\omega$, for all
  large $s$; for a bounded operator with numerical range in that sector,
  $\norm{\e^{wA_s}}\leq\e^{(\abs\omega+1)\abs w}$ on $\{\abs{\im w}\leq\re w\}$.  The
  exponentials $\e^{wA_s}$ are then Cauchy uniformly on bounded parts of that closed sector, by
  the complex form of the telescoping estimate of Theorem~\ref{thm:dynamics:hille-yosida}; the
  strong limit is the extension, holomorphic in $w$ at every vector as a locally uniform limit
  of holomorphic maps, and locally bounded, hence holomorphic in the operator norm by the
  last statement (Cauchy estimates for the first two derivatives and the mean value
  inequality).  The open sector $\Sigma_{\pi/4}$ lies in $\{\abs{\im w}\leq\re w\}$.  For the
  restriction, the map $z\mapsto(\id-P_W)T(z)x$, $x\in W$, vanishes on $[0,\infty)$ and is
  holomorphic on the connected open sector, so it vanishes there by the identity theorem.
\end{proof}

\begin{lemma}[The compact embedding]
  \label{lem:dynamics:compact-embedding}
  \lean{FastDynamo.isCompactOperator_lpDiagonal,
    FastDynamo.isCompactOperator_productH2Embedding}
  \leanok
  \uses{def:main:induction-operator}
  The embedding $\Hs^2(\T^3;\CC^3)\hookrightarrow\Lp^2(\T^3;\CC^3)$ is a compact operator.
\end{lemma}

This is the compact embedding the paper's proof of Lemma~A.1 uses.  Module
\texttt{FastDynamo/Operator/ProductH2Compact.lean}.

\begin{proof}
  \leanok
  In Fourier variables the embedding is division of the $\Hs^2$-weighted coefficients by the
  weight $1+4\pi^2\abs\xi^2$, which tends to infinity; a diagonal operator on $\ell^2$ with
  multipliers tending to zero is the operator-norm limit of its finite truncations, which have
  finite rank.
\end{proof}

\begin{lemma}[Immediate compactness]
  \label{lem:dynamics:immediate-compactness}
  \lean{FastDynamo.IsAnalyticSemigroup.exists_generator_comp,
    FastDynamo.IsAnalyticSemigroup.isCompactOperator}
  \leanok
  \uses{def:dynamics:analytic-semigroup, def:dynamics:generator,
    lem:dynamics:generator-basic}
  Let $S$ be an analytic semigroup on a Banach space with generator $G$.  For every $t>0$,
  $S(t)$ maps into $D(G)$ and $GS(t)$ is bounded.  If moreover $G$ has a compact resolvent,
  then $S(t)$ is compact for every $t>0$.
\end{lemma}

The paper cites [Engel--Nagel, Theorem~II.4.29] for immediate compactness; the Lean proof is
the direct factorisation below.  Module
\texttt{FastDynamo/Dynamics/AnalyticSemigroupGenerator.lean}.

\begin{proof}
  \leanok
  \uses{lem:dynamics:generator-basic}
  The holomorphic extension $T$ is differentiable at the real point $t$, and its complex
  derivative there is a bounded operator $D$; the real difference quotient of $S$ at $S(t)x$
  converges to $Dx$, so $S(t)x\in D(G)$ and $GS(t)=D$.  Then
  $S(t)=R(r)\bigl(rS(t)-GS(t)\bigr)$ is a compact operator composed with a bounded one.
\end{proof}

\begin{theorem}[Spectral mapping and the growth bound]
  \label{thm:dynamics:spectral-mapping}
  \lean{FastDynamo.StronglyContinuousSemigroup.exp_re_mul_le_norm,
    FastDynamo.StronglyContinuousSemigroup.operator_natCast,
    FastDynamo.StronglyContinuousSemigroup.cexp_mem_spectrum,
    FastDynamo.IsAnalyticSemigroup.exists_cexp_eq_of_mem_spectrum,
    FastDynamo.IsAnalyticSemigroup.spectrum_diff_zero,
    FastDynamo.IsAnalyticSemigroup.growthBound_eq_iSup}
  \leanok
  \uses{lem:dynamics:immediate-compactness, lem:dynamics:generator-basic,
    lem:dynamics:growth-bound-limit, def:dynamics:growth-bound}
  Let $S$ be an analytic semigroup on a complex Banach space whose generator $G$ has a compact
  resolvent.  Then for every $t>0$
  \[
    \sigma(S(t))\setminus\{0\}=\{\e^{t\lambda}:\lambda\in\sigma(G)\},
  \]
  and if $\sigma(G)\neq\emptyset$, then
  $\omega_0(S)=\sup\{\re\lambda:\lambda\in\sigma(G)\}$ and the supremum is attained at an
  eigenvalue of $G$, with an eigenvector.
\end{theorem}

The paper cites [Engel--Nagel, Corollaries~IV.3.11 and~IV.3.12] for the spectral mapping and
the spectral bound, and [Engel--Nagel, Corollary~V.3.2] for the attainment.  The Lean proof is a
direct point-spectrum argument.  Module \texttt{FastDynamo/Dynamics/SemigroupSpectralMapping.lean}.

\begin{proof}
  \leanok
  \uses{lem:dynamics:immediate-compactness, lem:dynamics:generator-basic,
    lem:dynamics:growth-bound-limit}
  Every point of $\sigma(G)$ is an eigenvalue (compact resolvent), and an eigenvector
  $Gx=\lambda x$ gives $S(t)x=\e^{t\lambda}x$, so $\e^{t\lambda}\in\sigma(S(t))$ and
  $\e^{t\re\lambda}\leq\norm{S(t)}$.  Conversely $S(t)$, $t>0$, is compact
  (Lemma~\ref{lem:dynamics:immediate-compactness}), so a non-zero $\mu\in\sigma(S(t))$ is an
  eigenvalue with finite-dimensional eigenspace; that eigenspace is invariant under the bounded
  operator $GS(t)$, which therefore has an eigenvector there, and it is an eigenvector of $G$
  whose eigenvalue $\lambda$ has $\e^{t\lambda}=\mu$.  For the growth bound: the eigenvector
  bound gives $\re\lambda\leq\omega_0(S)$; by Lemma~\ref{lem:dynamics:growth-bound-limit} and
  Gelfand's formula $\e^{\omega_0(S)}$ is the spectral radius of the compact operator $S(1)$,
  which is attained at a non-zero spectral point $\e^\lambda$, $\lambda\in\sigma(G)$, with
  $\re\lambda=\omega_0(S)$.
\end{proof}

\begin{lemma}[Mean conservation]
  \label{lem:dynamics:mean-conservation}
  \lean{FastDynamo.HasZeroMean, FastDynamo.productCoeffCLM_zero_zero_eq_integral,
    FastDynamo.hasZeroMean_productInductionAction,
    FastDynamo.exists_productSolenoidal_mean_eq_one,
    FastDynamo.one_le_norm_solenoidalInductionSemigroup,
    FastDynamo.gamma_eq_lim, FastDynamo.zero_le_gamma,
    FastDynamo.zero_mem_spectrum_generator_solenoidalInductionSemigroup}
  \leanok
  \uses{thm:dynamics:solenoidal-generator, lem:dynamics:growth-bound-limit,
    lem:dynamics:generator-basic}
  Let $u$ be transverse divergence-free and $\eps>0$.  For every divergence-free
  $B\in\Hs^2$, $\int_{\T^3}L_{\eps,u}B\dd x=0$; $\Lp^2_\sigma$ contains a constant field of norm
  one; hence $0\in\sigma(L_{\eps,u})$ and $\norm{\e^{tL_{\eps,u}}}\geq1$ for every $t\geq0$.
  Consequently $\gamma(u,\eps)$ is a real number $L$ with
  $0\leq L\leq M_u$ and $t^{-1}\log\norm{\e^{tL_{\eps,u}}}\to L$.
\end{lemma}

The first three facts are the paper's argument for $0\in\sigma(L_{\eps,u})$ in the proof of
Lemma~A.1; the paper does not state $\norm{\e^{tL_{\eps,u}}}\geq1$ or $\gamma\geq0$.  The mean
is the integral over the product torus with its probability Haar measure, which is Lebesgue
measure on $\T^3$ (\texttt{HasZeroMean}, \texttt{FastDynamo/Finish/ZeroMean.lean}).  Modules
\texttt{FastDynamo/Dynamics/Assembly/InductionMeanConservation.lean} and
\texttt{FastDynamo/Dynamics/Assembly/InductionSemigroupSpectrum.lean}.

\begin{proof}
  \leanok
  \uses{lem:dynamics:growth-bound-limit, lem:dynamics:generator-basic}
  The zeroth Fourier coefficient is the mean; that of $\eps\Delta B$ and of the weak curl
  vanishes.  So $L_{\eps,u}$ maps into the zero-mean fields and is not surjective onto
  $\Lp^2_\sigma$, which contains the constant field $c=(1,0,0)$: $0\in\sigma(L_{\eps,u})$.
  The mean of $\e^{tL_{\eps,u}}c$ has right derivative zero, so it stays $1$, and
  $\norm{\e^{tL_{\eps,u}}c}\geq1=\norm c$.  The minorant $1\leq\norm{\e^{tL_{\eps,u}}}$ makes
  the limit of Lemma~\ref{lem:dynamics:growth-bound-limit} finite and non-negative, and the
  Hille--Yosida bound gives $L\leq M_u$.
\end{proof}

\begin{theorem}[The spectral theory of the induction semigroup]
  \label{thm:dynamics:induction-spectrum}
  \lean{FastDynamo.isAnalyticSemigroup_solenoidalInductionSemigroup,
    FastDynamo.exists_isResolvent_generator_solenoidalInductionSemigroup,
    FastDynamo.isCompactOperator_of_isResolvent_generator_solenoidalInductionSemigroup,
    FastDynamo.isCompactOperator_solenoidalInductionSemigroup,
    FastDynamo.spectrum_solenoidalInductionSemigroup_diff_zero,
    FastDynamo.gamma_eq_iSup_re_spectrum,
    FastDynamo.exists_eigenvalue_re_eq_gamma}
  \leanok
  \uses{thm:dynamics:solenoidal-generator, lem:dynamics:sectorial,
    thm:dynamics:analytic-generation, lem:dynamics:compact-embedding,
    lem:dynamics:immediate-compactness, thm:dynamics:spectral-mapping,
    lem:dynamics:mean-conservation, prop:dynamics:product-resolvent}
  Let $u$ be transverse divergence-free and $\eps>0$.  The semigroup $\e^{tL_{\eps,u}}$ on
  $\Lp^2_\sigma$ is analytic of angle $\pi/4$; its generator $L_{\eps,u}$ has a non-empty
  resolvent set and a compact resolvent at every point of it; $\e^{tL_{\eps,u}}$ is compact
  for $t>0$; \eqref{eq:dynamics:compact-semigroup-spectral-mapping} holds; and
  $\gamma(u,\eps)=\sup\{\re\lambda:\lambda\in\sigma(L_{\eps,u})\}$, attained at an
  eigenvalue.
\end{theorem}

These are the instances, at the induction semigroup, of the abstract statements of this
section; Lemma~\ref{lem:semigroup-spectral-growth} collects them with the generator and the
limit.  Module \texttt{FastDynamo/Dynamics/Assembly/InductionSemigroupSpectrum.lean}.

\begin{proof}
  \leanok
  \uses{lem:dynamics:sectorial, thm:dynamics:analytic-generation,
    lem:dynamics:compact-embedding, lem:dynamics:immediate-compactness,
    thm:dynamics:spectral-mapping, lem:dynamics:mean-conservation}
  Analyticity: Theorem~\ref{thm:dynamics:analytic-generation} at the product Hille--Yosida ray
  with the sector condition of Lemma~\ref{lem:dynamics:sectorial}, then restriction to
  $\Lp^2_\sigma$.  Compact resolvent: at a real $r$ beyond the ray constant the resolvent of
  the generator is the restricted product resolvent
  (\texttt{exists\_isResolvent\_generator\_solenoidalInductionSemigroup\_eq}), which maps into
  $\Hs^2$ and is compact by Lemma~\ref{lem:dynamics:compact-embedding}; the resolvent identity
  transfers compactness to every point of the resolvent set.  The rest is
  Lemma~\ref{lem:dynamics:immediate-compactness} and
  Theorem~\ref{thm:dynamics:spectral-mapping}, the spectrum being non-empty by
  Lemma~\ref{lem:dynamics:mean-conservation}, and the growth bound of the product semigroup
  being that of its transport to the literal torus.
\end{proof}

\section{Solutions of the induction equation}
\label{sec:dynamics:solutions}

\begin{lemma}[The energy identity]
  \label{lem:dynamics:energy-identity}
  \lean{FastDynamo.productGradientEnergy,
    FastDynamo.inner_productH2Embedding_laplacian_eq_neg_gradientEnergy,
    FastDynamo.re_inner_productInductionTransverseTransport_eq_zero,
    FastDynamo.re_inner_productInductionAxialTransport_eq_zero,
    FastDynamo.re_inner_productInductionAction_eq,
    FastDynamo.hasDerivAt_half_norm_sq_of_isForwardStrongSolution,
    FastDynamo.hasDerivAt_half_norm_sq_of_isStrongSolution,
    FastDynamo.realStrongSolution_hasDerivAt_half_norm_sq}
  \leanok
  \uses{def:main:induction-operator, def:dynamics:semigroup-interface}
  Let $u$ be transverse divergence-free and $\eps\in\R$.  For every $B\in\Hs^2(\T^3;\CC^3)$,
  \[
    \re\inner{B}{\mathcal A_\eps B}=\re\inner{B}{(B\cdot\nabla)u}-\eps\norm{\nabla B}_2^2,
  \]
  and along every strong solution $B$ of $\partial_tB=\mathcal A_\eps B$ (at every positive
  time, for a forward strong solution)
  \[
    \frac12\frac{\dd}{\dd t}\norm{B(t)}_2^2
    =\re\inner{B(t)}{(B(t)\cdot\nabla)u}-\eps\norm{\nabla B(t)}_2^2 .
  \]
  In particular this holds along the real solution $B^\eps$ of Theorem~\ref{thm:main}.
\end{lemma}

This is the paper's \texttt{eq:energy-competition}, which the paper displays formally (``this
competition is expressed by the energy identity'').  The Lean identity is for complex fields with
the real part of the pairing; for a real field the pairing is the paper's
$\int_{\T^3}B\cdot(B\cdot\nabla)u\dd x$.  The proof uses no divergence condition on $B$:
both transport terms of $\mathcal A_\eps$ drop out because $\Div u=0$ and $u$ is independent of
$x_3$.  Off divergence-free fields, however, $\mathcal A_\eps$ is not $L_{\eps,u}$, so the
paper's identity is the case $\Div B=0$, where the auxiliary and the curl forms agree; along
solutions of the curl-form equation itself with divergence-free datum the identity is part of
Lemma~\ref{lem:dynamics:physical-uniqueness}.  Module
\texttt{FastDynamo/Dynamics/EnergyIdentity.lean}.

\begin{proof}
  \leanok
  \uses{def:main:induction-operator}
  In Fourier variables $\re\inner{B}{\eps\Delta B}=-\eps\norm{\nabla B}_2^2$; the transverse
  transport term is skew because $\Div_\perp u_\perp=0$, and the axial transport term
  $u_3\partial_3B$ is skew because $u_3$ does not depend on $x_3$.  Along a strong solution,
  $\frac{\dd}{\dd t}\norm{B}_2^2=2\re\inner{B}{\mathcal A_\eps B}$.
\end{proof}

\begin{lemma}[Uniqueness of the solution of the induction equation]
  \label{lem:dynamics:physical-uniqueness}
  \lean{FastDynamo.productHMinus2Synthesis,
    FastDynamo.exists_norm_stateResolvent_le_productL2ToHMinus2,
    FastDynamo.productHMinus2DivergenceFunctional,
    FastDynamo.productDivergenceFromH2_eq_zero_of_isForwardStrongSolution,
    FastDynamo.productPhysicalInductionOperator_forwardStrongSolutionsUnique,
    FastDynamo.isForwardStrongSolution_productPhysicalInductionOperator_domainOrbit,
    FastDynamo.productH2Embedding_eq_productInductionSemigroup_of_isForwardStrongSolution,
    FastDynamo.productH2Embedding_eq_solenoidalInductionSemigroup_of_isForwardStrongSolution,
    FastDynamo.hasDerivAt_half_norm_sq_of_isForwardStrongSolution_physical}
  \leanok
  \uses{def:finish:physical-operator, lem:dynamics:product-dissipative,
    thm:dynamics:product-semigroup, thm:dynamics:solenoidal-semigroup,
    lem:finish:physical-compatibility, lem:dynamics:energy-identity}
  Let $u$ be transverse divergence-free and $\eps>0$.  Two forward strong solutions
  $B_1,B_2:[0,\infty)\to\Hs^2(\T^3;\CC^3)$ of
  $\partial_tB=\curl(u\times B)+\eps\Delta B$, the time derivative taken in $\Hs^{-2}$, with
  $B_1(0)=B_2(0)$ agree for all $t\geq0$.  A forward strong solution with divergence-free
  initial value stays divergence-free, and it is the orbit
  $B(t)=\e^{tL_{\eps,u}}B(0)$ of the semigroup on $\Lp^2_\sigma$; conversely the orbit of
  every divergence-free $\Hs^2$ datum is a forward strong solution.  Along a forward strong
  solution with divergence-free initial value the energy identity
  \texttt{eq:energy-competition} holds at every $t>0$:
  $\frac12\frac{\dd}{\dd t}\norm{B(t)}_2^2
  =\re\inner{B(t)}{(B(t)\cdot\nabla)u}-\eps\norm{\nabla B(t)}_2^2$.
\end{lemma}

This is what the definite article in ``\emph{its} solution'' of Theorem~\ref{thm:main} asks
for.  The paper has no separate statement: Theorem~\ref{thm:main} asserts the norm law for
``its solution'', and its proof writes that solution down.  The solution class is the Lean
one: $\Hs^2$-valued curves on the product torus, with one-sided derivative at $t=0$ and
two-sided derivative at $t>0$ in the Fourier model of $\Hs^{-2}$, and no divergence condition
on the data; conclusions are for $t\geq0$.  It contains the class of clause~(16) of
Definition~\ref{def:main:physical-clauses}, whose curve is a strong solution at every real
time.  Modules \texttt{FastDynamo/Dynamics/Assembly/NegativeSobolevResolvent.lean},
\texttt{FastDynamo/Dynamics/Assembly/PhysicalInductionUniqueness.lean} and
\texttt{FastDynamo/Dynamics/Assembly/SemigroupOrbit.lean}.

\begin{proof}
  \leanok
  \uses{lem:dynamics:product-dissipative, lem:finish:physical-compatibility,
    lem:dynamics:energy-identity}
  The divergence of a forward strong solution satisfies, Fourier mode by Fourier mode, the
  heat equation $\dot c=-4\pi^2\eps\abs\xi^2c$, because the curl term has no divergence; so a
  divergence-free datum stays divergence-free.  On divergence-free fields the curl form is the
  $\Hs^{-2}$ image of $\mathcal A_\eps$ (Lemma~\ref{lem:finish:physical-compatibility}).  For the
  difference $w$ of two solutions, the resolvent $R(r)$ of $\mathcal A_\eps$ at a real $r>M_u$ is
  bounded in the $\Hs^{-2}$ norm (by the energy inequality and the skewness of transport), so
  $R(r)w(t)$ is an $\Lp^2$ forward strong solution of $\partial_t=\mathcal A_\eps$ with initial
  value $0$, hence $0$ by Lemma~\ref{lem:dynamics:product-dissipative}; $R(r)$ is injective, so
  $w=0$.  The same argument with the orbit of the realization of
  Theorem~\ref{thm:dynamics:product-semigroup} as second solution identifies a solution with
  divergence-free datum with the semigroup orbit.  That orbit is an $\Lp^2$ forward strong
  solution of $\partial_t=\mathcal A_\eps$, so Lemma~\ref{lem:dynamics:energy-identity} gives
  the energy identity along it, and the two curves agree on $[0,\infty)$.
\end{proof}
